% Títulos obligatorios preservados; el detalle del registro se entrega en T-16.
\chapter{Introducción a los Riesgos}
Synaptix gestionará los riesgos de implementación, marchas blancas y Operación del Puerto Deportivo Bahía Panitao, con horizonte de 56 meses. El plan protege personas, integridad de datos, continuidad, plazos y obligaciones de la oferta. Integra requisitos de SD3/T-12, decisiones de SD4/SD5, métodos de SD6, paquetes y reservas de SD7, controles de calidad y recepción de SD9 e innovaciones de SD13. El archivo independiente \ArchivoTDieciseis{} desarrolla los dieciséis riesgos generales y dieciocho de innovaciones, sus respuestas y trazabilidad; los registros de ejecución se producirán durante el contrato. \parencites[Formulario T-7, Subdocumento 8, p.~58; T-16, p.~63]{sd8-ba}[RT-19.04, p.~33]{sd8-bt}

\section{Plan de riesgos}
El procedimiento \texttt{RIESGO-8} fija cómo identificar, evaluar, tratar, comunicar y revisar los eventos que pueden impedir los resultados comprometidos. Su aplicación no permite aceptar como riesgo una excepción a un requisito obligatorio.

\subsection{Objetivo, alcance y enfoque de gestión}
El plan aplica ISO 31000:2018 mediante contexto y criterios explícitos, identificación por causa--evento--consecuencia, análisis, evaluación, tratamiento y seguimiento. La comunicación y el registro acompañan todo el ciclo. ISO 31000 aporta directrices; este plan concreta responsables, controles y productos de Synaptix y no declara una certificación obtenida. \parencite{sd8-iso31000}

El contexto de evaluación es la custodia de embarcaciones ajenas, la presencia de menores en el agua, el expendio de combustible y la trazabilidad de residuos peligrosos, con personal estacional y un CLIENTE sin área TI. Estas condiciones orientan los criterios de daño y la asignación de responsabilidades; no son eventos inciertos. El rechazo a cámaras sobre menores constituye una restricción de diseño: el riesgo es incumplirla o que falle una alternativa permitida. \parencite[caps.~1--3 y 19]{sd8-c07}

El horizonte cubre diseño y construcción E1 en meses 1--12, marcha blanca E1 en 13--15 y producción desde 16; construcción E2 en 13--18, marcha blanca E2 en 19--20 y producción desde 21; y los 36 meses de operación entre 21 y 56. La evaluación conserva los solapamientos y los hitos obligatorios. \parencite[arts.~15 y 17]{sd8-ba}

Las perspectivas de solución, desarrollo e implantación de T-22 se registran separadamente de las categorías RBS. La perspectiva indica qué resultado se amenaza: funcionamiento de la solución; construcción y entrega del proyecto; o transición, adopción y puesta en producción. La RBS clasifica la causa en proyecto, técnica, organización, seguridad u operación. Por ejemplo, R-03 es técnico y afecta solución e implantación; R-05 pertenece a proyecto y afecta desarrollo; R-04 es organizacional y afecta implantación. Un riesgo puede afectar varias perspectivas sin duplicar su registro ni su reserva. \parencite[Formulario T-22, Informe 2]{sd8-ba}

La evaluación cualitativa aplica las escalas de esta sección y compara obligación, gravedad, exposición y dependencia. El análisis cuantitativo de 8.2 calcula sensibilidades de duración y esfuerzo a partir de escenarios explícitos; no convierte puntuaciones ordinales en probabilidades estadísticas. El tratamiento de 8.3 compara prevención, mitigación y contingencia por su capacidad, plazo y resultado comprobable. Un objetivo residual conserva la exposición inicial hasta verificar la eficacia.

La línea de referencia comprende el alcance SD3/T-12, el despliegue híbrido y autonomía de SD4, la autoridad y conciliación de datos de SD5, las puertas de SD6/SD9, la disponibilidad y ventanas de SD7/T-15/T-18 y las dependencias de IN1--IN5 en SD13. Si falla la autonomía se activa R-02 y quedan afectados los flujos críticos; si falla la autoridad de escritura, R-03 afecta integridad y reconexión; si faltan agenda o acta funcional, R-04 impide el corte; si no hay capacidad efectiva, R-05 bloquea la selección; si la medición o el piloto no llegan a la ventana, R-15 afecta evidencia de innovación. Las habilitaciones externas EXT de SD7 siguen siendo dependencias por comprobar, no compromisos ya recibidos. Las variantes, fuentes, permisos, volúmenes de migración y disponibilidad se reevaluarán ante cualquier cambio; la revisión anterior no sustituye esos documentos actuales.

Se distingue riesgo de problema, supuesto y restricción. Una cancelación futura por marejada es un riesgo de plazo; el aviso emitido activa una suspensión obligatoria y una incidencia de planificación. El congelamiento conocido es una restricción que se incorpora al calendario, no una probabilidad que pueda aceptarse. La participación de la Contraparte en la recepción continua es un supuesto que debe comprobarse; el acta futura sigue siendo una puerta de corte. Una deficiencia comprobada se gestiona como acción correctiva, conservando el riesgo de recurrencia.

La Figura~\ref{fig:sd8-proceso} relaciona el tratamiento con la ejecución y aceptación.
\begin{figure}[H]
\centering
\begin{tikzpicture}[
 caja/.style={draw=SynLine,fill=SynSurface,rounded corners=1mm,align=left,text width=12.4cm,inner sep=3mm,font=\fontsize{9}{12}\selectfont},
 enlace/.style={-{Stealth},draw=SynBlue,thick}]
\node[caja] (a) {\textbf{Contexto y criterios}\quad Proyecto y titulares: obligación, flujo, dependencia, ventana y capacidad.};
\node[caja,below=4mm of a] (b) {\textbf{Identificación y evaluación}\quad Causa, evento y consecuencia; RBS, P/I/E, modo de falla y sensibilidad.};
\node[caja,below=4mm of b] (c) {\textbf{Decisión y tratamiento}\quad Responsable, disparador, plazo, recurso y reserva única; aprobación competente antes de ejecutar.};
\node[caja,below=4mm of c] (d) {\textbf{Comprobación y revisión}\quad Ensayo/evidencia, residual observado y cierre; nueva acción si falla eficacia.};
\draw[enlace] (a)--(b);
\draw[enlace] (b)--(c);
\draw[enlace] (c)--(d);
\draw[enlace] (d.east)--++(8mm,0)|-(b.east);
\end{tikzpicture}
\caption{Ciclo RIESGO-8 y control de eficacia}\label{fig:sd8-proceso}
\FuenteFigura{Elaboración de Synaptix a partir del enfoque ISO 31000:2018 y RT-19.04 de las Bases Técnicas Transversales, p.~33.}
\end{figure}

Proyecto registra la decisión y asigna capacidad antes de ejecutar una respuesta. El cierre exige evidencia de eficacia; cambiar una puntuación o emitir un acta de comité no demuestra que el control funcione. Una acción que modifique el compromiso contractual utiliza CC-6 de T-9 antes de ejecutarse.

La Tabla~\ref{tab:sd8-referencia} individualiza las condiciones que se comprueban antes de evaluar o autorizar trabajo. Son decisiones y dependencias de la oferta; la validación prevista no se presenta como efectuada.
\begingroup\setlength{\tabcolsep}{3pt}
\begin{longtable}{L{\dimexpr.25\textwidth-2\tabcolsep\relax}L{\dimexpr.37\textwidth-2\tabcolsep\relax}L{\dimexpr.38\textwidth-2\tabcolsep\relax}}
\caption{Línea de referencia y consecuencias de variación}\label{tab:sd8-referencia}\\\hline
\EncabezadoTabla{Condición y fuente} & \EncabezadoTabla{Validación y responsable} & \EncabezadoTabla{Riesgo y alcance afectado}\\\hline\endfirsthead
\hline\EncabezadoTabla{Condición y fuente} & \EncabezadoTabla{Validación y responsable} & \EncabezadoTabla{Riesgo y alcance afectado}\\\hline\endhead
Despliegue híbrido, autonomía y recuperación; SD4. & Arquitectura y SRE contrastan diseño y ensayos antes de producción y ante cambio. & R-02/R-07/R-08: custodia, zarpes, escuela y servicio; bloquear aceptación si falta una puerta.\\\hline
Autoridad, historia, volumen y conciliación; SD5 y DEC-02 de SD3. & Datos perfila el origen, conserva originales y comprueba retorno con Calidad antes del corte. & R-03/R-12: datos, saldos y expedientes; volumen o interpretación insuficientes obligan a recalcular y mantener convivencia.\\\hline
Recursos, agendas y tamaño funcional; T-15 vigente. & Proyecto contrasta carga/capacidad por persona, período y frente antes de seleccionar trabajo; UCP es sólo antecedente. & R-05/R-13/R-14: construcción, cobertura y adopción; falta de capacidad exige nivelación y relevo antes de ejecutar.\\\hline
Delegaciones funcionales, agenda y suplencia; DEC-22, SD3 y CASO 16.1. & Contraparte confirma autoridad y disponibilidad en arranque y antes de recepción; Proyecto conserva el acuerdo. & R-04/R-12: decisiones, cierre del legado y aceptación; ausencia bloquea la actividad dependiente, sin aceptación tácita.\\\hline
Provisión EXT y ventanas; SD7/T-14/T-18. & Proyecto y Arquitectura comprueban entrega, permisos e interfaces antes del consumidor; contrastan avisos y calendario. & R-01/R-06/R-11: montaje, pruebas e hitos; atraso exige propagación en la red y alternativa admisible.\\\hline
Medición, pilotos y dependencias IN1--IN5; SD13. & Proyecto y titulares comprueban criterios y evidencia antes del piloto y de la campaña dependiente. & R-15 y complemento de innovaciones: serie, adopción y beneficio; insuficiencia impide declarar resultado o madurez alcanzados.\\\hline
\end{longtable}\endgroup
\FuenteTabla{Síntesis de Synaptix a partir de SD3, SD4, SD5, SD7/T-15/T-18 y SD13; Caso 07, numerales 16.1 y 17.1.}
Cada variación conserva fuente, condición anterior, dato nuevo y efecto en riesgo, paquetes y puerta de aceptación. Proyecto abre la reevaluación del mismo ID y tramita control de cambios si corresponde; no usa un antecedente de esfuerzo ni una fecha ofrecida como línea base aprobada. Las designaciones y la capacidad se comprueban antes de usar la condición, sin ocultar su estado pendiente.

\subsection{Roles, responsabilidades y escalamiento}
Ignacio Silva, en Dirección de Proyecto, mantiene T-16, integra calendario y cuentas de costo, coordina decisiones y verifica reservas. Cada titular de riesgo responde por su tratamiento: Fernando Guerrero en Arquitectura, Bruno Toro en Seguridad, Davor Serey en Calidad y José Lara Arce en Operación/SRE. Fernando Guerrero coordina también Integración y Nicolás Torfan aporta evaluación y ejecución de datos y migración, conforme a la tabla de perfiles de T-15. La oferta asigna también a José Lara Arce la coordinación de Implantación desde el mes 8, alineada con T-15. Su función SRE comienza desde el mes 6. La dedicación, los ejecutores y los relevos se nivelan por período; la asignación conjunta no demuestra cobertura simultánea de esas tareas. Sus funciones compartidas consumen una sola capacidad por persona.

La Contraparte designa un referente y suplente por proceso; valida reglas y recepción sin asumir funciones de un área TI inexistente. El Comité de Proyecto revisa cada riesgo quincenalmente. El Comité Ejecutivo mensual resuelve exposición alta/crítica, aceptación residual competente y uso de reserva de gestión; el Comité de Arquitectura mensual revisa riesgos técnicos y deuda; el Comité de Operación mensual trata continuidad desde el mes 13. Las actas registran acuerdos, responsables y plazos dentro de los dos días hábiles siguientes a cada comité. \parencites[art.~71, pp.~37--38]{sd8-ba}[RT-19.04/09, pp.~33--34]{sd8-bt}[cap.~10, restricción 7, p.~26; numeral 13.3, pp.~29--30]{sd8-c07}

Ante peligro de personas, pérdida de integridad, incidente crítico o aviso de marejadas se activa de inmediato el procedimiento competente. La escala no retrasa esa acción. Seguridad y SRE conservan autoridades y notificaciones del incidente; Proyecto coordina recuperación de plazo y capacidad.

El propietario identifica y mantiene causa, exposición y respuesta; Proyecto coordina la evaluación con Arquitectura/Integración, Datos y Seguridad según el efecto; Calidad comprueba la evidencia de eficacia. Cada acción identifica ejecutor, competencia requerida, paquete, disponibilidad y fecha antes de autorizarse. El propietario puede ejecutarla si tiene competencia y agenda; coordinar el riesgo no le atribuye automáticamente todas las tareas técnicas. La asignación se contrasta con T-15, incluidos turnos, revisiones y ausencias, sin crear puestos adicionales no dimensionados.

La Contraparte Técnica coordina referentes funcionales de operaciones marítimas para amarras, custodia y zarpes; de escuela para presencia, autorizaciones y salud; de varadero para izaje, contratistas y residuos; de portería para ingreso y salida; y de administración para contratos, cobranza y cierre del legado. Son funciones del caso, no personal TI. El CLIENTE designará titular y suplente con autoridad funcional y agenda; los entrevistados no se consideran designados ni comprometidos por haber participado en el diagnóstico. Synaptix conserva administración técnica, accesos, respaldos y soporte. \parencite[numeral 2.4; numeral 16.1, decisión 22; numeral 17.6]{sd8-c07}

Ante ausencia de un titular, Proyecto transfiere expediente, permisos y tareas a un suplente de competencia equivalente y disponibilidad verificada; si no existe esa cobertura, escala R-04/R-05/R-13 y bloquea la decisión o actividad dependiente. Las identidades y agendas de suplentes técnicos y funcionales deberán confirmarse en el registro de recursos; no se afirma que estén disponibles. Para el sistema de 2014, Datos coordina el responsable de migración y su técnico suplente previstos en SD5, y la Contraparte designa un suplente funcional para cierre y conciliación. La documentación y transferencia temprana de SD3 se enlazan a R-12/R-04; antes del corte se exige repetir el procedimiento sin depender de la única persona conocedora del origen. Si falta esa evidencia, se mantiene convivencia y se bloquea el corte; el detalle de contingencia se integra con 8.2/8.3.

SRE dirige la activación de recuperación técnica conforme al procedimiento autorizado y Seguridad la respuesta de su ámbito; el referente funcional conserva las decisiones de seguridad física de su proceso. Proyecto coordina la contingencia de ejecución y comunica al CLIENTE estado, efecto, medidas y próxima decisión, sin esperar una reunión ante gravedad. La aceptación residual respeta las competencias de los comités y no sustituye la recepción del CLIENTE. Si la respuesta cambia la línea base se tramita CC-6 de T-9 con impacto en alcance, plazo, costo, riesgo y servicio, acuerdo escrito de ambas partes y aprobación del Ejecutivo antes de ejecutar. No se modifican objeto, naturaleza contractual, garantías mínimas ni meses obligatorios. \parencites[art.~72]{sd8-ba}[RT-19.03]{sd8-bt}

Para una ausencia prevista, Proyecto comprueba y registra el relevo antes de liberar la agenda del titular; ante ausencia imprevista suspende la tarea dependiente, protege evidencias y activa el escalamiento hasta habilitar cobertura equivalente. El acta de relevo identifica perfil, permisos, tareas, estado, pendientes, agenda y aceptación de la transferencia. La autoridad contractual o funcional se delega por el procedimiento del CLIENTE; no se transfiere sólo por disponer de acceso técnico.

La contingencia temprana del legado comienza en el mes 6 según SD3: inventario y copia recuperable, procedimientos de extracción y cierre, interpretación documentada y simulacro de conciliación. Datos coordina el ejecutor técnico y su suplente previstos en SD5; Administración designa el suplente funcional, que debe ejecutar el cierre documentado durante el segundo ensayo. Calidad comprueba repetibilidad y Seguridad los permisos. Si la persona conocedora del origen falta antes de completar esa transferencia, se preservan originales y último cierre comprobado, se mantiene el servicio de origen en la condición segura disponible y se bloquean extracción, interpretación o corte que carezcan de evidencia suficiente. Proyecto comunica el efecto y acuerda recuperación funcional sin inventar saldos ni presumir conocimiento transferido. Durante convivencia, la conciliación diaria y un cierre mensual completo permiten comprobar la sustitución; las discrepancias críticas bloquean el avance. Los riesgos R-04/R-12/R-13 conservan seguimiento hasta habilitar relevo y eficacia. \parencite[cap.~19; numeral 17.6]{sd8-c07}

\subsection{Criterios de evaluación y aceptación de riesgos}
La oferta utiliza escalas ordinales de uno a cinco. Las valoraciones iniciales son hipótesis de priorización basadas en exposición y diseño, sin frecuencia histórica ni probabilidad estadística acreditadas. La Tabla~\ref{tab:sd8-probabilidad} fija el significado de $P$.
\begin{table}[htbp]
\caption{Escala ordinal de probabilidad}\label{tab:sd8-probabilidad}
\begin{tabularx}{\linewidth}{L{12mm}L{29mm}Y}
\hline
\EncabezadoTabla{P} & \EncabezadoTabla{Nivel} & \EncabezadoTabla{Criterio de clasificación}\\\hline
1 & Excepcional & Requiere coincidencia de condiciones poco expuestas; justificar mecanismo y controles.\\\hline
2 & Poco probable & Exposición acotada y tratamiento previsto; demostrar eficacia antes de adoptar el residual.\\\hline
3 & Posible & Existe un mecanismo concreto y exposición; faltan datos o verificaciones para descartarlo.\\\hline
4 & Probable & Condición recurrente o varios frentes expuestos, como avisos y participación operacional limitada.\\\hline
5 & Muy probable & Desencadenante próximo o repetición comprobada; revisar si el evento ya es un problema.\\\hline
\end{tabularx}
\FuenteTabla{Criterios de priorización definidos por Synaptix para RIESGO-8.}
\end{table}
No se traduce $P=4$ a un porcentaje. El titular documentará motivo, exposición, incertidumbre y fuente; nuevos ensayos, incidentes o cambios pueden modificar la clasificación sin borrar su versión anterior.

El impacto $I$ toma el mayor efecto entre plazo, continuidad, calidad de datos, privacidad, personas, custodia, obligaciones ambientales y aceptación. La Tabla~\ref{tab:sd8-impacto} expresa umbrales de oferta, subordinados a las obligaciones contractuales.
\begin{table}[htbp]
\caption{Escala de impacto y exposición}\label{tab:sd8-impacto}
\begin{tabularx}{\linewidth}{L{12mm}L{42mm}Y}
\hline
\EncabezadoTabla{I} & \EncabezadoTabla{Efecto de plazo} & \EncabezadoTabla{Efecto operacional o de aceptación}\\\hline
1 & Hasta un día hábil & Corrección local sin afectar personas, datos ni criterios obligatorios.\\\hline
2 & Hasta dos días hábiles & Retrabajo acotado, recuperable dentro de la capacidad y ventana disponible.\\\hline
3 & Hasta cinco días hábiles & Varios resultados requieren recuperación; degradación reversible no crítica.\\\hline
4 & Hasta diez días hábiles & Se compromete margen de varios frentes o una preparación de recepción.\\\hline
5 & Más de diez días o hito afectado & Peligro de personas, pérdida de datos, exposición grave o incumplimiento de un mínimo obligatorio.\\\hline
\end{tabularx}
\FuenteTabla{Criterios de Synaptix; los límites contractuales prevalecen sobre esta escala.}
\end{table}
Con $E=P I$, los intervalos son bajo 1--4, moderado 5--9, apreciable 10--14, alto 15--19 y crítico 20--25. $E$ es un índice ordinal, no pérdida monetaria esperada. Un impacto 5 en personas, custodia, integridad de registros, privacidad u obligaciones ambientales exige escalamiento inmediato y revisión competente aunque el producto sea menor a 15. El incumplimiento ambiental grave o la pérdida de trazabilidad de custodia se clasifica por ese efecto, no sólo por días de atraso; una baja probabilidad no autoriza rebajar la protección ni aceptar una condición impeditiva.

El titular puede proponer aceptación baja, que Proyecto registra; el Comité de Proyecto decide moderada/apreciable dentro de su competencia. Alto/crítico se eleva al Ejecutivo con respuesta, capacidad y fecha. Ninguna aceptación autoriza un incumplimiento, falta de prueba impeditiva o corte sin acta. El riesgo residual se valora sólo tras comprobar controles; T-16 muestra un objetivo residual diferenciado del residual observado, todavía no acreditado.

\subsection{Registro y trazabilidad de los riesgos}
El registro versionado \path{componentes/registro_riesgos.csv} contiene R-01 a R-16 y es la fuente de las fichas T-16. Cada fila conserva categoría RBS, causa, evento, consecuencia, $P/I/E$ iniciales, objetivo residual, titular, disparador, mitigación, contingencia, plazo, evidencia, paquetes, reserva, fuente y estado. Una nueva revisión añade evaluación residual observada, fecha, evaluador competente y referencia de eficacia.

La Tabla~\ref{tab:sd8-enlace-pl} relaciona los cinco riesgos de planificación existentes sin renombrarlos ni reemplazar sus responsables.
\begin{table}[htbp]
\caption{Integración con los riesgos de SD7}\label{tab:sd8-enlace-pl}
\begin{tabularx}{\linewidth}{L{23mm}L{27mm}Y}
\hline
\EncabezadoTabla{SD7} & \EncabezadoTabla{T-16} & \EncabezadoTabla{Tratamiento integrado}\\\hline
PL-R01 & R-01 & Suspensión de todo trabajo programado; controlar capacidad global y reserva marina.\\\hline
PL-R02 & R-03/R-12 & Arquitectura conserva responsabilidad del retrabajo; Calidad verifica la evidencia afectada.\\\hline
PL-R03 & R-12 & Repetición de ensayo y esfuerzo de corrección separados de las dos ejecuciones obligatorias.\\\hline
PL-R04 & R-04 & Recepción continua condicionada a participación y acta expresa.\\\hline
PL-R05 & R-05 & Sensibilidad UCP, capacidad real y refuerzo propio.\\\hline
\end{tabularx}
\FuenteTabla{Integración de T-16 con \path{riesgos_reservas_planificacion.csv} y las fichas de SD7.}
\end{table}
Un mismo evento puede relacionarse con varias fichas, pero conserva una sola incidencia, respuesta, consumo y cuenta de costo. Los IDs T-12 y complementos del caso permanecen íntegros; el tratamiento no elimina una cláusula ni cambia su estado automáticamente.

\subsection{Consolidación de innovaciones y registro complementario}\label{sec:sd8-innovaciones}
Las dieciséis fichas generales se conservan. El complemento \path{componentes/registro_riesgos_innovaciones.csv}, desarrollado dentro del T-16 en la sección ``Consolidación de los riesgos de IN1--IN5'', aplica la misma escala a dieciocho mecanismos de SD13. Reúne adopción, integración y acceso de IN1 (R-IN1-01--04); cohorte, atrasos, marejadas y dependencia de IN1 para invernada de IN2 (R-IN2-01--04); extracción, revisión humana, tratamiento autorizado, continuidad y beneficio de IN3 (R-IN3-01--05); atribución y omisión de casos difíciles por el incentivo de IN4 (R-IN4-01/02); y validez ambiental, participación y privacidad entre flotas de IN5 (R-IN5-01--03).

Cada mecanismo mantiene causa, consecuencia, valoración inicial, objetivo residual, titular, disparador, mitigación, contingencia, hito, evidencia, paquetes y estado. El vínculo con R-01/02/03/10/14/15 permite gestionar dependencias sin sumar puntuaciones ordinales ni duplicar una reserva. Los responsables de control coordinan las funciones de SD13; no reemplazan las atribuciones de sus autores, revisores o receptores. Proyecto revisa ambos registros en cada Comité de Proyecto y ante el disparador. Una omisión de solicitudes difíciles, un cierre no sustentado o un deterioro de calidad en IN4 suspende el reconocimiento variable y exige conciliar el universo completo, incluidas abiertas y reaperturas, manteniendo servicio fijo y SLA.

Los residuales son objetivos aún no observados; sólo evidencia de eficacia y decisión competente permiten adoptarlos. Los hitos conservan el calendario de SD13/SD7, la incorporación operativa del mes 21 y la ventana estacional IN2 de 2029. No se afirma evaluación de beneficio ni madurez ya alcanzada; el registro atiende la obligación general RT-19.04 aunque un criterio de innovación sea deseable. Las acciones y refuerzos se nivelan en T-15 y se valorizan exclusivamente en OE, sin reserva monetaria creada por esta consolidación. \parencite[RT-19.04, p.~33]{sd8-bt}

\subsection{Ciclo de revisión, seguimiento y comunicación}
Proyecto incorpora el riesgo en el arranque y lo revisa en cada Comité de Proyecto, además de cambios arquitectónicos, fallas de piloto, resultados de migración, cambios de ventana, incidentes graves, indisponibilidad de actores críticos, avisos y cercanía de ventanas críticas. Estos eventos obligan a reevaluar sin esperar la reunión siguiente. El titular revisa el disparador al seleccionar trabajo, antes de promover una versión y al preparar una campaña. Publica estado, acción próxima y evidencia en el espacio colaborativo de T-9 accesible al CLIENTE, con registro de riesgos y cambios siempre actualizados y datos sensibles restringidos por perfil. Cada actualización conserva fecha, versión, causa, evaluación inicial y observada, acción, ejecutor, plazo, exposición, estado y evidencia; no borra el antecedente.

Los estados son identificado, evaluado, tratamiento autorizado, en ejecución, en seguimiento, materializado y cerrado. Un riesgo materializado se enlaza al incidente y mantiene seguimiento del residual. Para cerrar, el titular aporta evidencia, Calidad comprueba eficacia y Proyecto registra decisión y lección; Seguridad/Operación revisan su ámbito. El informe mensual de avance incluye riesgos, desviaciones, incidencias y compromisos del período siguiente, junto con exposición y consumo de reservas. Si cambia la causa, falla un control o aparece nueva evidencia tras el cierre, Proyecto reabre el mismo ID, enlaza el incidente o cambio y solicita una evaluación y respuesta nuevas. Desviación mayor a 10\,\% en un hito exige aviso en cinco días hábiles y recuperación, sin esperar al comité siguiente. \parencite[RT-19.04--10, pp.~33--34]{sd8-bt}

\section{Identificación y Análisis de Riesgos}
La RBS distingue riesgos de proyecto, técnicos, organizacionales, de seguridad y de operación. La identificación parte de flujos y dependencias concretos: autonomía local, integración con legado, personas menores, estructura flotante, turnos estacionales, dependencia del proveedor y recepción de las etapas. T-16 reúne causas y consecuencias; este capítulo analiza prioridad y sensibilidad.

\subsection{RBS y prioridad de tratamiento}
La Tabla~\ref{tab:sd8-rbs} agrupa los dieciséis riesgos por mecanismo y criterio de comprobación.
\begin{table}[htbp]
\caption{RBS y exposición inicial de la oferta}\label{tab:sd8-rbs}
\begin{tabularx}{\linewidth}{L{27mm}L{42mm}L{17mm}Y}
\hline
\EncabezadoTabla{Grupo} & \EncabezadoTabla{Riesgos} & \EncabezadoTabla{E máx.} & \EncabezadoTabla{Comprobación decisiva}\\\hline
Proyecto & R-01/05/12/15/16 & 20 & Ventanas/capacidad, dos ensayos, serie/pilotos y depósito por liberación.\\\hline
Técnico & R-03/06/07/08 & 15 & Autoridad/conciliación, variante, salida y carga representativa.\\\hline
Organizacional & R-04/14 & 20 & Referente/suplente, recepción expresa y competencia individual.\\\hline
Seguridad & R-09/10 & 15 & Procedencia, intrusión y rechazo de acceso indebido.\\\hline
Operación & R-02/11/13 & 15 & Autonomía/DR real, seguridad física y cobertura efectiva.\\\hline
\end{tabularx}
\FuenteTabla{Evaluaciones iniciales de \path{registro_riesgos.csv}; puntuaciones ordinales de oferta.}
\end{table}
R-01 y R-04 alcanzan 20 por suspensión y aceptación sin holgura; requieren atención antes de cada ventana, no sólo respuesta posterior. R-14 alcanza 16 por renovación estacional y capacitación protegida. Otros seis riesgos tienen 15 y siete tienen 12. La prioridad de una falla de permisos o integridad permanece alta por su impacto aunque una valoración residual futura reduzca su exposición.

El orden combina obligación, dependencia y exposición. No se posterga ciberseguridad por disponer de una puntuación menor a la de plazo. Obsolescencia y bloqueo por proveedor se controlan desde la selección y el plan de salida; el peak se comprueba con carga, no con el promedio anual. La disponibilidad de la Contraparte requiere confirmación por proceso y no puede reemplazarse con un supuesto de aceptación tácita. \parencites[cap.~10, pp.~26--27; numerales 13.1--13.3, pp.~28--30; cap.~15, pp.~33--36]{sd8-c07}[RT-19.04, p.~33; cap.~20, pp.~34--35]{sd8-bt}

\subsection{Análisis de modos de falla}
Se selecciona análisis de modos de falla para ordenar las verificaciones de diseño y recepción, con IEC 31010 como referencia para la selección de técnicas. Además de $P$ e $I$, se utiliza detectabilidad $D$, criterio ordinal propuesto por Synaptix: 1, comprobación directa antes del efecto; 2, comprobación documental previa que permite detectar la ausencia de una puerta; 3, ensayo o conciliación dedicada; 4, ensayo negativo de fallas o permisos cuya omisión puede ocultar el daño; 5, detección tardía o evidencia insuficiente. Cada valor describe el mecanismo de detección, sin atribuir eficacia ya probada; la Tabla~\ref{tab:sd8-fmea} muestra hipótesis iniciales. \parencite{sd8-31010}
\begin{table}[htbp]
\caption{Modos de falla prioritarios y comprobación}\label{tab:sd8-fmea}
\begin{tabularx}{\linewidth}{Y L{23mm}L{20mm}Y}
\hline
\EncabezadoTabla{Modo de falla} & \EncabezadoTabla{P/I/D} & \EncabezadoTabla{NPR} & \EncabezadoTabla{Control y resultado exigido}\\\hline
R-03: hecho duplicado al reconectar & 3/5/4 & 60 & Fallos/reintentos y conciliación sin pérdida ni duplicado; sincronización en plazo.\\\hline
R-10: permiso expone datos de menores & 3/5/4 & 60 & Casos negativos por rol/finalidad; denegación y auditoría correctas.\\\hline
R-12: retorno no representa hechos nuevos & 3/5/3 & 45 & Ensayar restauración y reaplicación con volumen y respuestas pendientes.\\\hline
R-04: acta ausente en el corte & 4/5/2 & 40 & Agenda confirmada y evidencia completa; bloquear corte sin firma expresa.\\\hline
\end{tabularx}
\FuenteTabla{Análisis de Synaptix: $NPR=P I D$, con supuestos ordinales y controles SD4/SD5/SD7.}
\end{table}
NPR no es probabilidad ni monto esperado y no se utiliza para autorizar una puerta incumplida. La mayor prioridad de detección en datos y permisos justifica ensayos negativos y conciliación. Tras el ensayo se revisa $D$ con evidencia, no por la sola existencia del control.

\subsection{Sensibilidad de plazo, esfuerzo y exposición}
La suspensión por marejadas afecta toda actividad programada. Para mostrar capacidad expuesta se usa un escenario de ocho ejecutores con cinco HH directas por día, la referencia de SD7. La Tabla~\ref{tab:sd8-marejada} multiplica duración suspendida por esa capacidad; no estima frecuencia meteorológica ni afirma pérdida real.
\begin{table}[htbp]
\caption{Escenario de capacidad suspendida}\label{tab:sd8-marejada}
\begin{tabularx}{\linewidth}{L{38mm}L{38mm}Y}
\hline
\EncabezadoTabla{Suspensión} & \EncabezadoTabla{HH expuestas} & \EncabezadoTabla{Decisión requerida}\\\hline
5 días hábiles & $5\times8\times5=200$ & Comprobar reprogramación y reserva por todos los frentes.\\\hline
10 días hábiles & $10\times8\times5=400$ & No asumir que diez días marinos restituyen estas horas.\\\hline
15 días hábiles & $15\times8\times5=600$ & Supera el escenario marino; dimensionar refuerzo y red antes del hito.\\\hline
\end{tabularx}
\FuenteTabla{Escenario de Synaptix; capacidad directa de referencia SD7, sin convertir ocho titulares en disponibilidad efectiva simultánea.}
\end{table}
La operación de un servicio existente y la emergencia conservan su procedimiento; no se utilizan para continuar trabajo de proyecto suspendido. Con dotación distinta se calcula $H_{\mathrm{exp}}=\sum_{d,p} h_{p,d}^{\mathrm{programado}}$ sobre agenda y alcance bloqueados. Las HH expuestas no son horas adicionales automáticas: el esfuerzo de recuperación depende de reprogramación, revisión, ventanas y recursos habilitados.

PL-R02 utiliza $(d,d,1{,}6d)$ para retrabajo, con duración PERT $1{,}1d$ y desviación de escenario $0{,}1d$; no aporta percentiles ni presupuesto de corrección. PL-R03 usa $(10,10,16)$ días: $(10+4\times10+16)/6=11$ y $(16-10)/6=1$ día. La subred de SD7 distingue los déficits de ocho días E1 y dos E2 del escenario esperado, que no se compensan con renombrar una puerta.

Para R-05 se utiliza el antecedente funcional que conserva el T-15 vigente: 5.257,66 HH con factor UCP 20 y 7.360,72 HH con factor 28. La diferencia es 2.103,06 HH; $(7.360{,}72-5.257{,}66)/5.257{,}66\simeq40\,\%$. El cambio de factor supone el mismo tamaño funcional antecedente y no equivale a un aumento de 40\,\% de toda la EDT. No se adopta como presupuesto seleccionado ni se convierte en una reserva disponible. Para trabajo vigente se calcula por actividad $\Delta H=H_{\mathrm{reestimado}}-H_{\mathrm{seleccionado}}$ y se contrasta la carga por perfil con agenda y capacidad de T-15. Sin esfuerzo seleccionado y recursos suficientes, el resultado contractual permanece pendiente de cuantificación. Proyecto bloquea trabajo sobreasignado y nivela o asigna refuerzo comprobado antes de ejecutarlo.

Los análisis de R-07/R-08 conservan la estimación de salida y los peaks declarados en SD4; una transferencia contractual no equivale a reconstruir otra plataforma, y las pruebas a 1,5 veces peak no prueban toda demanda futura. Los compromisos de autonomía, RTO/RPO y tiempos de negocio permanecen puertas verificables, no riesgos aceptables por puntuación.

\subsection{Fundamento de valoración y dependencia entre riesgos}
La probabilidad ordinal describe la posibilidad de materialización durante los 56 meses del contrato, en las actividades y ventanas expuestas de cada ficha. No es una tasa mensual ni una frecuencia por maniobra. Los valores 2, 3 y 4 se fundamentan en condición limitada, mecanismo plausible sin historial suficiente y exposición recurrente o concentrada, respectivamente. La Tabla~\ref{tab:sd8-fundamentos} individualiza la causa y el criterio empleado; todos son juicios iniciales de ingeniería, sujetos a nueva evidencia. El impacto toma la consecuencia máxima de cada ficha, con escalamiento independiente de la multiplicación cuando afecta la protección obligatoria.
\begingroup\setlength{\tabcolsep}{3pt}
\begin{longtable}{L{.10\textwidth}L{.27\textwidth}L{.27\textwidth}L{.27\textwidth}}
\caption{Fundamento de las valoraciones iniciales}\label{tab:sd8-fundamentos}\\\hline
\EncabezadoTabla{ID} & \EncabezadoTabla{Condición concreta} & \EncabezadoTabla{Probabilidad propuesta} & \EncabezadoTabla{Impacto propuesto}\\\hline\endfirsthead
\hline\EncabezadoTabla{ID} & \EncabezadoTabla{Condición concreta} & \EncabezadoTabla{Probabilidad propuesta} & \EncabezadoTabla{Impacto propuesto}\\\hline\endhead
R-01 & Avisos de la autoridad marítima y ventanas costeras limitadas. & $P=4$: Hay recurrencia estacional, dependencia humana o ventana obligatoria que concentra la exposición. & $I=5$: Puede impedir un hito o afectar personas, integridad, privacidad o condición obligatoria; escalamiento por gravedad.\\\hline
R-02 & Dependencias físicas y fallos correlacionados en el recinto. & $P=3$: Existe causa y consecuencia identificables; sin historial suficiente para descartar el evento. & $I=5$: Puede impedir un hito o afectar personas, integridad, privacidad o condición obligatoria; escalamiento por gravedad.\\\hline
R-03 & Reintentos, doble escritura o reglas cloud/local incompatibles. & $P=3$: Existe causa y consecuencia identificables; sin historial suficiente para descartar el evento. & $I=5$: Puede impedir un hito o afectar personas, integridad, privacidad o condición obligatoria; escalamiento por gravedad.\\\hline
R-04 & Cliente sin TI, turnos operacionales y validación final sin holgura. & $P=4$: Hay recurrencia estacional, dependencia humana o ventana obligatoria que concentra la exposición. & $I=5$: Puede impedir un hito o afectar personas, integridad, privacidad o condición obligatoria; escalamiento por gravedad.\\\hline
R-05 & Funciones coincidentes, refuerzos no habilitados o factores de ambiente distintos del modelo. & $P=3$: Existe causa y consecuencia identificables; sin historial suficiente para descartar el evento. & $I=4$: Se comprometen habilitación, cobertura, ventana o calidad de varios resultados.\\\hline
R-06 & Fin de soporte, ficha incorrecta o componente no apto para el entorno marino. & $P=3$: Existe causa y consecuencia identificables; sin historial suficiente para descartar el evento. & $I=4$: Se comprometen habilitación, cobertura, ventana o calidad de varios resultados.\\\hline
R-07 & Servicios administrados, formatos o claves no exportables. & $P=3$: Existe causa y consecuencia identificables; sin historial suficiente para descartar el evento. & $I=4$: Se comprometen habilitación, cobertura, ventana o calidad de varios resultados.\\\hline
R-08 & Concentración operacional y frecuencia de medición mayor a la dimensionada. & $P=3$: Existe causa y consecuencia identificables; sin historial suficiente para descartar el evento. & $I=4$: Se comprometen habilitación, cobertura, ventana o calidad de varios resultados.\\\hline
R-09 & Origen no verificado, vulnerabilidad, secreto expuesto o procedencia falsificable. & $P=3$: Existe causa y consecuencia identificables; sin historial suficiente para descartar el evento. & $I=5$: Puede impedir un hito o afectar personas, integridad, privacidad o condición obligatoria; escalamiento por gravedad.\\\hline
R-10 & Permisos amplios, registros inseguros o datos reales en pruebas. & $P=3$: Existe causa y consecuencia identificables; sin historial suficiente para descartar el evento. & $I=5$: Puede impedir un hito o afectar personas, integridad, privacidad o condición obligatoria; escalamiento por gravedad.\\\hline
R-11 & Flexión, agua, combustible clasificado o interacción durante maniobra. & $P=3$: Existe causa y consecuencia identificables; sin historial suficiente para descartar el evento. & $I=5$: Puede impedir un hito o afectar personas, integridad, privacidad o condición obligatoria; escalamiento por gravedad.\\\hline
R-12 & Legado sin soporte/documentación y conocimiento concentrado. & $P=3$: Existe causa y consecuencia identificables; sin historial suficiente para descartar el evento. & $I=5$: Puede impedir un hito o afectar personas, integridad, privacidad o condición obligatoria; escalamiento por gravedad.\\\hline
R-13 & Turnos, suplencias, inglés o especialistas sin provisión efectiva. & $P=3$: Existe causa y consecuencia identificables; sin historial suficiente para descartar el evento. & $I=4$: Se comprometen habilitación, cobertura, ventana o calidad de varios resultados.\\\hline
R-14 & Personal nuevo y monitores renovados, con capacitación restringida en peak. & $P=4$: Hay recurrencia estacional, dependencia humana o ventana obligatoria que concentra la exposición. & $I=4$: Se comprometen habilitación, cobertura, ventana o calidad de varios resultados.\\\hline
R-15 & Dependencias de las innovaciones y ciclos naturales; medición tardía. & $P=3$: Existe causa y consecuencia identificables; sin historial suficiente para descartar el evento. & $I=4$: Se comprometen habilitación, cobertura, ventana o calidad de varios resultados.\\\hline
R-16 & Depósito por liberación omitido, titularidad incorrecta o custodia sin habilitar. & $P=3$: Existe causa y consecuencia identificables; sin historial suficiente para descartar el evento. & $I=4$: Se comprometen habilitación, cobertura, ventana o calidad de varios resultados.\\\hline
R-IN1-01 & Se conserva la coordinación telefónica sin ingresar sus cambios. & $P=4$: Hay recurrencia estacional, dependencia humana o ventana obligatoria que concentra la exposición. & $I=4$: Se comprometen habilitación, cobertura, ventana o calidad de varios resultados.\\\hline
R-IN1-02 & Adhesión voluntaria y pocas solicitudes elegibles. & $P=3$: Existe causa y consecuencia identificables; sin historial suficiente para descartar el evento. & $I=3$: Retrabajo o pérdida reversible de beneficio; no autoriza rebajar seguridad ni una puerta obligatoria.\\\hline
R-IN1-03 & Contratos, identidades o versiones difieren entre IN1, CTR y VAR. & $P=3$: Existe causa y consecuencia identificables; sin historial suficiente para descartar el evento. & $I=4$: Se comprometen habilitación, cobertura, ventana o calidad de varios resultados.\\\hline
R-IN1-04 & Autorización insuficiente de armador o contratista asignado. & $P=2$: Exposición condicionada a una autorización o vínculo específico; no consta fallo reiterado. & $I=5$: Puede impedir un hito o afectar personas, integridad, privacidad o condición obligatoria; escalamiento por gravedad.\\\hline
R-IN2-01 & Adhesión voluntaria o convocatoria sin grupo estable. & $P=3$: Existe causa y consecuencia identificables; sin historial suficiente para descartar el evento. & $I=4$: Se comprometen habilitación, cobertura, ventana o calidad de varios resultados.\\\hline
R-IN2-02 & Trabajo o disponibilidad no confirmado antes de la ventana. & $P=4$: Hay recurrencia estacional, dependencia humana o ventana obligatoria que concentra la exposición. & $I=4$: Se comprometen habilitación, cobertura, ventana o calidad de varios resultados.\\\hline
R-IN2-03 & Aviso marítimo o condición insegura coincidente con trabajos programados. & $P=4$: Hay recurrencia estacional, dependencia humana o ventana obligatoria que concentra la exposición. & $I=5$: Puede impedir un hito o afectar personas, integridad, privacidad o condición obligatoria; escalamiento por gravedad.\\\hline
R-IN2-04 & La coordinación preventiva depende de un servicio o contrato aún no recibido. & $P=3$: Existe causa y consecuencia identificables; sin historial suficiente para descartar el evento. & $I=4$: Se comprometen habilitación, cobertura, ventana o calidad de varios resultados.\\\hline
R-IN3-01 & Documento ilegible o formato cambiado; extractor propone valores erróneos. & $P=4$: Hay recurrencia estacional, dependencia humana o ventana obligatoria que concentra la exposición. & $I=4$: Se comprometen habilitación, cobertura, ventana o calidad de varios resultados.\\\hline
R-IN3-02 & Sesgo de automatización, formación incompleta o revisión apresurada. & $P=3$: Existe causa y consecuencia identificables; sin historial suficiente para descartar el evento. & $I=4$: Se comprometen habilitación, cobertura, ventana o calidad de varios resultados.\\\hline
R-IN3-03 & Proveedor, ubicación, retención o permiso difieren de condiciones aprobadas. & $P=3$: Existe causa y consecuencia identificables; sin historial suficiente para descartar el evento. & $I=5$: Puede impedir un hito o afectar personas, integridad, privacidad o condición obligatoria; escalamiento por gravedad.\\\hline
R-IN3-04 & Dependencia remota no disponible durante ingreso. & $P=3$: Existe causa y consecuencia identificables; sin historial suficiente para descartar el evento. & $I=3$: Retrabajo o pérdida reversible de beneficio; no autoriza rebajar seguridad ni una puerta obligatoria.\\\hline
R-IN3-05 & Volumen bajo o revisión consume el ahorro de captura. & $P=3$: Existe causa y consecuencia identificables; sin historial suficiente para descartar el evento. & $I=3$: Retrabajo o pérdida reversible de beneficio; no autoriza rebajar seguridad ni una puerta obligatoria.\\\hline
R-IN4-01 & Mezcla, dotación o mejoras ajenas cambian respecto de la base; medición no aprobada. & $P=3$: Existe causa y consecuencia identificables; sin historial suficiente para descartar el evento. & $I=5$: Puede impedir un hito o afectar personas, integridad, privacidad o condición obligatoria; escalamiento por gravedad.\\\hline
R-IN4-02 & Selección de tareas fáciles o presión para aparentar menor esfuerzo. & $P=3$: Existe causa y consecuencia identificables; sin historial suficiente para descartar el evento. & $I=5$: Puede impedir un hito o afectar personas, integridad, privacidad o condición obligatoria; escalamiento por gravedad.\\\hline
R-IN5-01 & Adhesión voluntaria, lecturas faltantes o cambios de ocupación, clima y lavado. & $P=3$: Existe causa y consecuencia identificables; sin historial suficiente para descartar el evento. & $I=4$: Se comprometen habilitación, cobertura, ventana o calidad de varios resultados.\\\hline
R-IN5-02 & Plan poco comprensible o abandono voluntario de armador/charter. & $P=3$: Existe causa y consecuencia identificables; sin historial suficiente para descartar el evento. & $I=3$: Retrabajo o pérdida reversible de beneficio; no autoriza rebajar seguridad ni una puerta obligatoria.\\\hline
R-IN5-03 & Vínculo histórico/consentimiento o autorización por embarcación no se comprueba. & $P=2$: Exposición condicionada a una autorización o vínculo específico; no consta fallo reiterado. & $I=5$: Puede impedir un hito o afectar personas, integridad, privacidad o condición obligatoria; escalamiento por gravedad.\\\hline
\end{longtable}\endgroup
\FuenteTabla{Registros generales y de innovaciones T-16; criterios de Synaptix, no frecuencias medidas.}
Los riesgos no son independientes: R-01/R-06/R-11 y R-IN2-02/03 pueden consumir la misma ventana física; R-03/R-12 comparten autoridad y conciliación; R-04/R-14 y riesgos de adopción comparten agenda funcional; R-02/R-13 comparten recuperación y cobertura. Proyecto agrupa el incidente, propaga a los consumidores y cuenta una sola suspensión o consumo. La valoración no suma productos ordinales ni multiplica probabilidades no disponibles. IN3 mantiene revisión humana aunque se estime ahorro; IN4 no mejora exposición ocultando casos abiertos o difíciles.

\subsection{Escenarios cuantificados y límites de recuperación}
Las magnitudes siguientes complementan FMEA. Se distingue presupuesto de error mensual, respuesta de incidente y RTO: cumplir cuatro horas en DR no permite gastar cuatro horas todos los meses bajo una disponibilidad crítica de 99,9\,\%. Para un mes de treinta días, $30\times24\times60\times(1-0{,}999)=43{,}2$ minutos de indisponibilidad crítica máxima; un incidente de cuatro horas excedería ese presupuesto mensual en 196,8 minutos. La medición se realiza sobre transacciones de negocio, con denominador y exclusiones contractuales, no sobre el estado del servidor. \parencite[arts.~20 y 78]{sd8-ba}

El escenario de restauración desde cero de SD5 utiliza 60 GB históricos provisionales, 6,25 MB/s de transferencia, 20 MB/s de restauración estructurada y 65 minutos de preparación, reaplicación y comprobación integrada. Se obtiene $65+60.000/(6{,}25\times60)+10.000/(20\times60)=233{,}33$ minutos; con doble volumen, 401,67 minutos. A volumen doble, conservar 6,25 MB/s no cumple cuatro horas: se requieren al menos $120.000/[60(240-65-20.000/(20\times60))]=12{,}63$ MB/s efectivos de transferencia, manteniendo los otros supuestos. Este límite calculado no acredita una red ni elimina crecimiento, tiempos de identidad o fallas de copias; se ensaya y se ajusta capacidad antes de aceptar. La conmutación desde réplica preparada utiliza otra secuencia y no se confunde con restauración desde cero.

Para mostrador, 90 legajos a 90 segundos representan 8.100 segundos, 135 minutos de trabajo serial. Cuatro ejecutores sin otras tareas reducirían el lote ideal a 33,75 minutos, pero no demostrarían la latencia individual de 90 segundos ante llegadas simultáneas ni su disponibilidad real. El análisis obliga a preparar antecedentes locales, medir cola y tiempo completo y dimensionar puestos con la demanda y distribución de llegadas; el promedio anual no prueba el peak. Si falta capacidad o autorización externa, el sistema conserva el legajo y la tarea pendiente, sin conceder zarpe. \parencite[RT-09.01; numeral 16.1, decisión 21]{sd8-c07}

La holgura de marejada es el escenario de diez días hábiles por campaña registrado en SD7, no un resultado meteorológico ni una holgura CPM demostrada. Si el margen admisible antes del hito es $M$ y la suspensión consume $s$, el déficit de calendario es $\max(0,s-M)$. Con $M=10$ días y suspensiones de 5/10/15, quedan 5/0/0 días y aparece déficit 0/0/5; son escenarios condicionados a disponer realmente de esos diez días fuera del congelamiento. Se retira del programa toda actividad suspendida, incluida remota; dos ventanas alternativas no se suman por existir en el documento. La red vigente de SD7/T-15 aún tiene duraciones y asignaciones pendientes: el plan identifica esa brecha, y no afirma cabida ni cumplimiento de hitos por restar escenarios aislados. Si el margen real es cero, cualquier suspensión positiva requiere recuperación demostrada o registra atraso, sin mover meses contractuales.

El perfil preparado de SD4 conserva 18 vCPU/72 GiB frente al peak de 32/128: el mínimo aporta sólo 56,25\,\% del perfil peak y requiere agregar 14 vCPU/56 GiB antes de admitir ese tráfico. La espera de ampliación pertenece al RTO. Mantener el mínimo es una reducción de capacidad permanente de 43,75\,\% frente a duplicar el peak, no ahorro monetario demostrado. Prueba, cuota y acceso a capacidad antes de H5/H10 controlan R-02/R-08; una promoción sin expansión y prueba de carga no se acepta.

\subsection{Cobertura de identificación y dependencias excluidas}\label{sec:sd8-cobertura-identificacion}
La identificación recorre alcance y restricciones SD3/T-12, arquitectura e inventario SD4/T-11, integración y datos SD5, métodos SD6, EDT y calendarios SD7/T-14/T-15/T-18, pruebas SD9 e innovaciones SD13. El registro \path{componentes/trazabilidad_riesgos.csv} individualiza para los 34 IDs categoría causal, perspectivas T-22, fuente, paquete, horizonte y referencia al análisis. Conserva los IDs, propietarios y valores de las fichas; los 18 riesgos de innovaciones son escenarios propios vinculados a un riesgo general, no copias adicionales de la misma reserva.

La revisión del alcance detecta fallas de aceptación y capacidad en R-04/05/12; la arquitectura, obsolescencia, portabilidad y saturación en R-02/03/06/07/08; seguridad de entrega y acceso en R-09/10; ejecución física y cobertura en R-01/11/13; formación, medición y cierre en R-14/15/16. Los complementos IN1--IN5 distinguen integración, adopción, exposición documental, incentivo y atribución ambiental. La RBS clasifica estas causas en cinco grupos; cada categoría del complemento figura también en la matriz oficial de T-16. No se fuerza una fila por requisito ni se inventa una cantidad mínima de riesgos. \parencites[Formulario T-7, apartado 10; T-16; T-22]{sd8-ba}[cap.~8]{sd8-ex}

La exclusión del sistema contable no elimina el riesgo de respuesta incierta de pago o factura: R-03/R-12 requieren consulta por ID y conciliación antes de repetir un efecto. La autoridad marítima conserva la autorización de zarpe; R-02/R-03 consideran su falta de respuesta y conservan el legajo pendiente sin conceder autorización. R-06/R-11 cubren entrega de hardware, licencias y habilitación de obras/recinto del CLIENTE mediante EXT-EQUIPOS/LICENCIAS/RECINTO; R-07/R-13 cubren dependencia de proveedor y soporte. No se incluyen esos sistemas como desarrollos propios ni se supone su recepción. \parencite[cap.~11; numeral 17.1]{sd8-c07}

Cuando un evento produce efectos diferentes se conserva una respuesta por efecto: R-11 exige detener la maniobra insegura; R-10 aislar acceso a datos protegidos; R-01 suspender y reprogramar; R-02 recuperar continuidad; R-03/R-12 conciliar; R-15 mantener evidencia ambiental y no declarar certificación obtenida. R-IN4-02 protege calidad frente al incentivo, y R-IN5-01 separa atribución ambiental de continuidad de participación R-IN5-02. Un evento compartido puede activar varias respuestas, pero su suspensión y reserva se consumen una sola vez. Los 1.150 movimientos anuales de travelift son exposición operacional del caso, no 1.150 daños ni una probabilidad de falla.

\subsection{Detalle reproducible de los modos de falla}\label{sec:sd8-fmea-detalle}
Se utiliza FMEA porque permite seguir una función concreta hasta el efecto, el control y la prueba sin disponer de tasas estadísticas de falla. No se seleccionan árbol de fallas ni simulación: faltan probabilidades de eventos básicos, independencia y distribuciones verificables; no se calculan percentiles ficticios. Los cálculos físicos de sensibilidad complementan FMEA, cuyo NPR permanece ordinal. La Tabla~\ref{tab:sd8-fmea-detalle} desarrolla la función, causa y efecto de los cuatro modos ya puntuados en 8.2.
\begingroup\setlength{\tabcolsep}{3pt}
\begin{longtable}{L{.11\textwidth}L{.26\textwidth}L{.28\textwidth}L{.25\textwidth}}
\caption{Función, causa y consecuencia de los modos prioritarios}\label{tab:sd8-fmea-detalle}\\\hline
\EncabezadoTabla{ID} & \EncabezadoTabla{Función y falla} & \EncabezadoTabla{Causa y efecto} & \EncabezadoTabla{Cálculo y decisión}\\\hline\endfirsthead
\hline\EncabezadoTabla{ID} & \EncabezadoTabla{Función y falla} & \EncabezadoTabla{Causa y efecto} & \EncabezadoTabla{Cálculo y decisión}\\\hline\endhead
R-03 & Confirmar un hecho una sola vez; duplicarlo al reconectar. & Reintento, identidad o autoridad no conciliada; duplica movimiento, despacho o efecto administrativo. & $3\times5\times4=60$; probar pérdida/duplicación y bloquear efecto no conciliado.\\\hline
R-10 & Consultar sólo con finalidad y facultad; exponer salud de menor. & Permiso o vínculo incorrecto; divulgación y decisión no autorizadas. & $3\times5\times4=60$; ensayar denegación por rol/finalidad y aislar el acceso fallido.\\\hline
R-12 & Retornar sin perder hechos posteriores; volver a una imagen anterior incompleta. & Restauración sin inventario/reaplicación; descarta pagos, despachos o respuestas pendientes. & $3\times5\times3=45$; dos ensayos y conciliación de hechos posteriores antes de cortar.\\\hline
R-04 & Obtener recepción expresa antes de habilitar; cortar sin acta. & Referente ausente, agenda o expediente incompleto; activación sin recepción y atraso. & $4\times5\times2=40$; confirmar agenda/suplente y bloquear corte sin firma.\\\hline
\end{longtable}\endgroup
\FuenteTabla{Causas y efectos de las fichas T-16; escalas P/I/D declaradas por Synaptix y controles SD4/SD5/SD7.}
Detectabilidad 4 en R-03/R-10 exige prueba negativa específica, 3 en R-12 ensayo de conciliación y 2 en R-04 comprobación documental de la puerta. El orden 60/60/45/40 orienta preparación de evidencia, sin rebajar el escalamiento por impacto 5 ni demostrar eficacia ejecutada.

\subsection{Contraste de ventanas y rutas con T-18}\label{sec:sd8-contraste-ventanas}
T-18 incorpora una red global de 2.279 nodos y 5.077 enlaces, con inicio propuesto el 1 de diciembre de 2026. Su tabla de programación presenta márgenes hábiles hacia obligaciones: H3 2,33 días, H5 11, H10 14,67, corte E1 1,44 y acta/corte E2 cero. Las cuatro ramas críticas de IN1/IN3/IN4/IN5 convergen en PILOTOS-REV/PILOTOS-SUB y H12; holgura cero hacia cierre de implementación no equivale a margen hacia un hito intermedio. Estas magnitudes proceden de \path{componentes/red_global.tex} de SD7 y \path{componentes/cpm_implementacion.csv} de SD7. T-15 conserva cálculos y cargas pendientes y otro antecedente funcional: la presencia de una red en T-18 no elimina esa discrepancia ni demuestra recursos para las 34 respuestas.

La reserva marina de diez días por campaña ya imputada por T-18 no se resta otra vez a sus márgenes. Como contraste, si una suspensión total $s$ excede la duración $r$ realmente incorporada en la rama, la prolongación adicional es $a=\max(0,s-r)$ y el déficit hacia su obligación es $\max(0,a-M)$, donde $M$ es el margen de esa misma rama. Con $r=10$, $s=15$ y $M=2{,}33$, el exceso es cinco días y el déficit 2,67; con $M=11$ el déficit aritmético es cero y quedan seis días de margen. Sólo valen esas cifras si reserva, actividad y margen pertenecen a la misma cadena y no cambian ruta, ventana o recursos. No se usa margen de H5 para justificar un atraso de H3 ni para financiar el corte E2.

Si una suspensión no estaba incluida en el frente, $r=0$: la demora adicional consume su margen completo. Una demora positiva en la firma final E2, cuyo margen es cero, incumple la puerta; no la absorben diez días marinos situados antes de marcha blanca. Como la marejada suspende todos los frentes programados, las ramas concurrentes reciben el mismo bloqueo de calendario y se recalculan juntas, sin sumar cinco días por cada rama como si fueran sucesivas. Si luego comparten ejecutores, la nivelación puede aumentar la demora respecto de ese cálculo de precedencias.

Los déficits de ocho días E1/dos E2 del contraste PERT de la subred antigua son resultados de ese escenario, no fechas finales de la red global de T-18. El antecedente funcional de T-15 5.257,66/7.360,72 HH conserva su frontera; la memoria \path{componentes/estimacion_funcional.tex} de SD7 plantea una selección distinta de 12.708 HH y sensibilidad 19.640,50 HH, con otro conteo, reparto y factores. No se sustituyen silenciosamente ni se suman ambos modelos: Proyecto debe resolver la referencia conjunta y propagar la selección a cargas y red. Tampoco se presenta el incremento entre modelos como la sensibilidad de 40\,\% del antecedente de T-15.

Esta sensibilidad identifica qué hitos quedan expuestos y por qué; no constituye simulación global ni cabida acreditada de toda cancelación. Para cerrar esa cabida se requieren fechas y duraciones de cada actividad afectada, calendario protegido, frentes, recursos y reserva realmente insertada, con nueva propagación de precedencias y nivelación. Hasta conciliar T-15/T-18 se conserva la brecha y se bloquea el compromiso que dependa de ella. Los retrasos calculados son adversos condicionados; no son días esperados porque no se dispone de una probabilidad estadística del evento. \parencites[art.~17]{sd8-ba}[cap.~10, restricción 12; numeral 17.5; cap.~19]{sd8-c07}

\subsection{Escenarios técnicos, de datos y seguridad}\label{sec:sd8-tecnicos}
El registro \path{componentes/escenarios_tecnicos_riesgos.csv} relaciona cada escenario con los IDs ya evaluados en T-16, sus valores iniciales, causa, efecto, tratamiento, contingencia y referencia técnica. Amplía la cobertura sin crear otra reserva ni cambiar una valoración por suponer que el control ya funciona. Los riesgos R-02/03/06/07/08/09/10/11/12/13 y los de IN3 concentran esta revisión. Un efecto confirmado sobre personas, privacidad, integridad o un mínimo obligatorio activa escalamiento por gravedad y reevaluación aunque la ficha original tenga impacto menor.

\subsubsection{Conectividad, borde y dependencias comunes}
R-02 considera caída del enlace principal, del respaldo y de ambos por ducto, alimentación, llegada o equipamiento compartidos. SD4 selecciona fibra Entel y radio Gtd en DA-18; dos nombres comerciales no prueban diversidad física. Arquitectura verifica rutas y conmutación con cada enlace retirado y ambos fuera de servicio; la pérdida de WAN activa las autonomías locales de 24/12/8 horas, mientras pérdida de nodo, disco, LAN o energía exige recuperación local distinta. El par PostgreSQL síncrono/Pacemaker/Corosync con fencing BMC y RAID 1 no protege contra inundación, incendio, red común, fallo de ambos nodos o exclusión fallida. Se conserva la captura durable y se bloquea promoción sin excluir el escritor; perder ambos nodos no se presenta como simple desconexión de Internet. \parencites[art.~16.4]{sd8-ba}[RT-02.11; RT-03.10--17]{sd8-bt}

R-02/R-11 incluyen los cinco pantalanes y varadero: sombra radioeléctrica, marea, oleaje, flexión y corrosión pueden impedir registrar una entrega, despacho o movimiento. Fibra marina, alivio de tensión, gabinetes protegidos, WLAN segmentada y terminales offline de SD4 4.2 deben comprobar cobertura con marea completa y trabajo real, sin prometer cobertura en navegación. Ante pérdida de campo se conserva el hecho con hora de origen y estado pendiente, se usa coordinación VHF y se suspende la actuación que necesite una condición de seguridad no comprobable; un registro manual alternativo no acredita captura electrónica cumplida.

Se identifican dependencias físicas y humanas comunes: recinto y acceso técnico, generación/climatización, agua que depende de la bomba eléctrica, circuito de llegada, autoridad de escritura, custodio de claves, firma de contraparte y conocimiento del legado. R-02/R-04/R-11/R-12/R-13 impiden tratarlas como resueltas por duplicar servidores. SD4 propone módulo técnico junto a administración, circuitos A/B, dos UPS y climatización N+1; su dimensionamiento y recepción no equivalen al gabinete actual ni sustituyen comprobar autonomía, ambiente y acceso. SRE administra el sitio; el CLIENTE conserva coordinación funcional sin recibir un cargo TI inexistente.

DR-4/REG-REP-4 y el DRP de 8.3 analizan São Paulo/México Central: fallo zonal, regional, réplica atrasada, IAM/DNS/KMS/despliegue o proveedor común. Se mide corte SQL/objetos y no sólo un contador de retraso. R-06/R-07 exigen verificar servicio, variante, cuota y condiciones en ambas regiones antes de habilitar; una región futura o una ampliación aún no recibida no es una contingencia operativa. Portabilidad SQL/Parquet/S3/OpenAPI/OIDC no hace portables de forma automática permisos Cognito, claves KMS, reglas EventBridge/SNS/SQS ni servicios de operación AWS. La salida exige reconfigurar adaptadores, exportar y verificar datos/identidades, recifrar con claves nuevas, reconstruir infraestructura y probar negocio; conserva la estimación propia de SD4 y no inventa un plazo universal.

R-06 incorpora el término de soporte antes del mes 56. SD4/T-11 mantienen versiones y términos públicos, incluido el de EMQX Enterprise 6.3 LTS; Synaptix debe disponer de soporte/renovación o migración antes del término. La actualización puede romper esquema, adaptador, dispositivo o réplica y coincidir con congelamiento: se ensaya regresión/retorno y se programa ventana admisible antes de obsolescencia, sin trasladar el riesgo al CLIENTE ni usar esa ventana para prorrogar remediación obligatoria.

\subsubsection{Capacidad, mensajes e integración}
R-08/R-13 evalúan la mañana con hasta 90 zarpes, visitantes y clases simultáneas, no el promedio anual. APIs, escritura PostgreSQL, radio, colas, proveedores de avisos y personas pueden limitar el servicio en puntos diferentes. SD4 declara perfiles cloud/local, prioridades y presupuestos de reconexión; SD5 conserva mezclas y escenarios distintos de carga. Sin una prueba conjunta no puede afirmarse qué recurso satura primero: Calidad medirá cola, latencia de negocio, CPU/I/O, enlace y atención por la misma mezcla, y bloqueará la aceptación si alguno incumple. Los 135 minutos seriales/33,75 ideales a cuatro puestos calculados en 8.2 no acreditan puestos disponibles ni latencia individual.

Lecturas por minuto, agregación de quince minutos, cola local y ACK durable de SD4 generan ráfaga al reconectar. R-03/R-08 consideran contención con analítica y avisos; sólo se elimina cola tras acuse íntegro, se prioriza captura crítica y se limita carga diferible. El crecimiento a tres veces volumen inicial en tres años incluye documentos, histórico, auditoría e índices; no es la ampliación evaluada a 380 amarras/70 boyas ni sólo triplicar sesiones sobre datos vacíos. Ambas sensibilidades se prueban separadamente con sus recursos y conservación. \parencites[RT-05.05; RT-09.01--09]{sd8-bt}[numerales 13.2 y 14.1/14.2]{sd8-c07}

R-03/R-12 cubren contabilidad/pagos, autoridad marítima, meteorología/VHF, avisos, firma y residuos. El catálogo A4.2 y los contratos de SD5 deben identificar canal realmente recibido, timeout y resultado. Caída o lentitud deja una tarea durable visible; error o formato cambiado conserva original en cuarentena y no publica un estado válido; repetición usa clave idempotente; resultado externo incierto se consulta por ID y se concilia antes de repetir cobro, despacho o aviso. Sin canal confirmado se usa procedimiento humano documentado, sin inventar API oficial ni declarar sincronización automática mediante transcripción. Firma no válida bloquea el documento que la necesita; falta de acuse genera escalamiento humano y no se registra como entrega comprobada. Aviso meteorológico oficial activa suspensión aunque falle el proveedor de notificación.

Los mensajes voluntarios de navegación pueden llegar tarde, desordenados, repetidos o no llegar. Se conservan hora del hecho, recepción, secuencia, procedencia y calidad; un silencio no confirma regreso, una posición antigua no es ubicación actual y un reintento no cierra otra salida. R-03/R-10 mantienen expediente/alerta y coordinación humana sin exigir seguimiento. La reconexión tampoco permite sobrescribir asistencia, perder despacho, duplicar hecho facturable o cerrar deuda con una respuesta no conciliada. SD4 conserva autoridad por período y SD5 inbox/efecto/outbox y AUD-5: conflicto queda visible y bloquea el efecto dependiente hasta resolución autorizada.

\subsubsection{Privacidad, conservación y prueba}
R-12 evalúa extracción incompleta, duplicados y cambio de identidad, armador o saldo del sistema de 2014. Originales, correspondencias, dos ensayos, cierres conciliados y retorno con hechos posteriores de SD5 evitan interpretar una copia como saldo confirmado; la ausencia del conocedor activa la contingencia del legado, sin inventar saldos. R-10 cubre salud/menores, deuda, contactos, tripulantes, contratistas y posición voluntaria. Un rol estacional, cuenta compartida o acceso de tercero puede ampliar visibilidad indebidamente: autorización de recurso/finalidad, mínimo privilegio, cifrado de campo, AUD-5 y autoridad local de SD4 limitan acceso, sin ocultar la brecha de decisiones externas durante aislamiento total.

Consentimiento inexistente, no comprobable o revocado impide habilitar posición voluntaria; una autorización general del dispositivo no lo reemplaza. Al reconectar se aplica revocación antes de reabrir tratamiento y no se impone seguimiento como solución a la falta de consentimiento. R-03/R-10/R-12 cubren pérdida, corrupción, divulgación y falta de evidencia como efectos distintos: autorizaciones, avisos/acuses, planes de izaje, actas ambientales, estado de cuenta y expedientes de las 14 embarcaciones abandonadas requieren procedencia, versiones y reconstrucción, además de una copia recuperable.

AUD-5/5.2.5 de SD5 rigen retención, bloqueo y eliminación de originales, auditoría, nube, borde, derivados y respaldos. Una categoría sin fecha de inicio o regla validada no se purga automáticamente; una restauración reaplica revocaciones/eliminaciones antes de servir y no reduce conservación probatoria a rotación de respaldo. SEC-ROT-1.0/GLS-1.0 y custodia independiente controlan pérdida de clave, certificado vencido y acceso de emergencia: se prueba descifrado, vigencia y doble custodia, sin destruir claves de evidencias aún exigibles ni abrir tráfico sin acceso autorizado. R-07/R-10 e IN3-03 incluyen transferencias a Brasil/México, soporte y subencargados; no se presupone permiso legal por la ubicación de AWS. Seguridad verifica tratamiento e instrucciones del CLIENTE antes de cualquier flujo. \parencites[arts.~21, 23 y 85]{sd8-ba}[RT-05.03/07; RT-07.10; RT-11.08--10]{sd8-bt}

\subsubsection{Amenazas y desarrollo seguro}
SEC-1.0 de SD4, sección «Modelo de amenazas y decisiones de seguridad», y A4.1 aplican STRIDE por componente y frontera. R-09/R-10/R-02/R-13 relacionan ransomware con custodia/datos, DDoS y abuso API con consulta/aviso/zarpe, privilegio indebido con entrega escolar y deuda, secretos expuestos con integración y cifrado, y dispositivo comprometido con captura/autoridad local. Cada cambio de región, versión mayor, permiso o integración reabre ese modelo y el riesgo; no se copia una aprobación anterior al activo nuevo.

El salto desde invitados/socios o CCTV hacia operación amenaza autenticidad de un despacho o aviso. VLAN/ACL y rutas separadas de SD4 deben rechazar acceso lateral, aun con credencial de segmento comprometida. R-09 incorpora dependencia vulnerable o artefacto alterado: PRUEBAS-6/SEC-1.0, SAST/SCA/DAST, escaneo, SBOM, firma/procedencia y DEP-6 impiden promover sin control. Desarrollo/preproducción usan datos sintéticos o protegidos, sin copiar producción indiscriminadamente ni habilitar acceso directo de desarrolladores. Se rechaza archivo no examinado; el nombre de WAF o del escáner no acredita cobertura completa.

La observabilidad de SD4/SD5 reúne OpenTelemetry/CloudWatch/CloudTrail, XDR/SIEM, eventos de negocio y registros protegidos. Se ensayan pérdida de colector, reloj, enlace, alteración de log, alerta sin respuesta y ausencia del analista. CloudTrail no sustituye AUD-5 ni una alerta prueba atención. El pool SOC propuesto para un puesto activo 24\,$\times$\,7 conserva relevo y escalamiento, sin presentarse como dotación contratada; R-13 mantiene capacidad pendiente en T-15. La pérdida de consola central conserva protección/alertas locales y exige respuesta autorizada, evitando aislar a la vez ambos nodos y perder custodia.

\subsubsection{IA documental y amenazas propias de IN3}\label{sec:sd8-ia-riesgos}
SD4 declara Textract y SD13 IN3 OCR/extracción de formularios; SD1 gobierna su uso bajo NIST AI RMF 1.0 e ISO/IEC 42001. IN5 y ocupación usan reglas deterministas, sin extenderles artificialmente el bloque IA. IN3-01/02 evalúan sesgo por formato/calidad y propuestas sin respaldo —alucinación o extracción incorrecta— que podrían incorporar una vigencia, monto o embarcación erróneos. IN3-03 evalúa fuga por envío, permisos o tratamiento y uso indebido de documentos o finalidad. NIST AI RMF se aplica mediante gobierno (inventario/responsables), contexto (documentos/campos/usuarios), medición (errores por campo/formato y resultados adversos) y gestión (bloqueo, revisión, desactivación y reevaluación). ISO/IEC 42001 aporta control versionado de propósito, cambios, supervisión y responsabilidades; no se atribuye certificación al extractor.

El modelado específico recorre original/cuarentena S3, adaptador, servicio Textract, salida propuesta, revisor y dato aprobado. Suplantación de revisor o documento exige identidad/finalidad; alteración de original/salida exige huella, versión y vínculo; repudio exige entrada/salida/corrección/decisión en AUD-5; divulgación exige permisos y tratamiento autorizado; denegación o resultado tardío exige ingreso manual; elevación de privilegio exige que el adaptador no escriba maestros aprobados. La confirmación humana no convierte una salida sin fuente en hecho cierto. IN3-C01--C05 y paquete IN3.4.2 de SD13 vinculan pruebas de rechazo de escritura, reconstrucción, permisos, regresión y desactivación; el cambio de esta frontera reabre SEC-1.0 conforme RT-26.07.

Antes de habilitar se fijan API/configuración/versionado aplicable, región efectiva y condiciones de proveedor que garanticen no entrenamiento de terceros salvo autorización expresa escrita. SD4 mantiene esa comprobación pendiente: no se infiere región del endpoint ni se envían datos reales para completar el piloto. Entrada, salida, confianza cuando la provea, fragmento original y decisión humana se conservan con acceso y retención de la categoría; la interfaz identifica propuesta automática. Si fallan calidad, permiso o servicio, SRE desactiva el adaptador, Datos/Calidad revisan pendientes por ID y el ingreso manual conserva original y revisión sin sobrescribir confirmados. No hay modelo predictivo propio ni plan de reentrenamiento automático seleccionado; cambios de formato/servicio se detectan con métricas y revalidación antes de habilitar. \parencites[art.~86]{sd8-ba}[RT-18.01--09; RT-26.07]{sd8-bt}


\subsection{Continuidad: procesos y dependencias afectadas}\label{sec:sd8-contexto-continuidad}
R-02/R-03/R-12/R-13 enlazan indisponibilidad, integridad, migración y capacidad de servicio con custodia, salidas y regresos, presencia de menores, combustible, residuos y administración. Su análisis de impacto de negocio debe distinguir seguridad de personas, interrupción admisible, pérdida admisible de datos y dependencia de energía, enlace, plataforma local y referente funcional. SD4 fija la arquitectura y los objetivos de continuidad; SD5 fija originales, conciliación y retorno. Esas decisiones son entradas del BCP/DRP conforme a ISO 22301 e ISO/IEC 27031 y no acreditan por sí solas su ejecución ni sustituyen los procedimientos de negocio. \parencites[Formulario T-7, apartado 10]{sd8-ba}[RT-10.03/04]{sd8-bt}

\subsection{BCP: impacto de negocio y niveles de servicio}\label{sec:sd8-bcp-bia}
El BCP aplica el ciclo de ISO 22301: identificar procesos y dependencias, analizar impacto, seleccionar continuidad, ejecutar una respuesta autorizada, comprobar recuperación y mejorar con resultados. El DRP aplica ISO/IEC 27031 a los recursos TIC de esa misma cadena. No se declara certificación ni simulacro realizado. Proyecto conserva versiones; SRE mantiene procedimientos y Seguridad custodia accesos; la Contraparte valida condiciones funcionales y prioridades. \parencites[Formulario T-7, apartado 10; arts.~20 y 78]{sd8-ba}[RT-10.02--04]{sd8-bt}

La Tabla~\ref{tab:sd8-bia} clasifica por consecuencia, no por tecnología. Las funciones críticas incluyen captura y consulta que permiten demostrar custodia y seguridad; una salida administrativa o un informe diferible puede pertenecer a otra clase dentro del mismo dominio. Para servicios críticos se comprometen RTO máximo de cuatro horas y RPO regional máximo de quince minutos. En captura desconectada, el requisito es conservar íntegramente los hechos críticos: ese RPO no permite descartar quince minutos de hechos locales. Las autonomías del caso son límites de capacidad de operación, no autorización para continuar una maniobra insegura.
\begingroup\setlength{\tabcolsep}{3pt}
\begin{longtable}{L{.22\textwidth}L{.10\textwidth}L{.29\textwidth}L{.29\textwidth}}
\caption{Impacto de negocio, dependencia y modo seguro}\label{tab:sd8-bia}\\\hline
\EncabezadoTabla{Servicio} & \EncabezadoTabla{Clase} & \EncabezadoTabla{Impacto/dependencia} & \EncabezadoTabla{Continuidad y límite}\\\hline\endfirsthead
\hline\EncabezadoTabla{Servicio} & \EncabezadoTabla{Clase} & \EncabezadoTabla{Impacto/dependencia} & \EncabezadoTabla{Continuidad y límite}\\\hline\endhead
Custodia, amarras, salida/regreso y legajo de zarpe & Crítica & Entrega indebida o desconocer quién está en el agua; autoridad local, diario y radio. & Recinto 24 h sin enlace; legajo local y coordinación humana; sin autorización no se concede zarpe ni entrega.\\\hline
Escuela: presencia, retirada, salud y emergencias & Crítica & Menor sin ubicación o entrega no autorizada; lista vigente, permisos y monitor. & Recinto 24 h y terminal 12 h; comprobación humana inmediata; si no hay certeza se suspende nueva salida y se aplica emergencia.\\\hline
Despacho de combustible y seguridad de maniobra & Crítica & Daño físico y pérdida de trazabilidad; habilitación, control local y condiciones seguras. & Captura local 24 h; detener despacho/izaje si falla la condición de seguridad o el registro exigible.\\\hline
Aviso urgente y contacto de afectados & Crítica & Falta de instrucción a personas expuestas; destinatarios, acuse, radio/telefonía alternativa. & Registro de llamada por no acuse; no presumir entrega ni cobertura marítima continua.\\\hline
Gestión de faena, acreditación y entrega de residuos & Alta & Se impide faena o entrega sin habilitación/trazabilidad; manifiesto, gestor y permisos. & Captura de terreno 12 h; detener operación dependiente sin autorización; conservar entrega/recepción probatorias.\\\hline
Administración: cierre, cobros y contabilidad & Alta & Diferencia de saldo o falta de expediente; legado, suplente funcional e interfaz externa. & Consulta local del último estado con antigüedad visible; diferir efectos externos y conciliar; no inferir saldo vigente.\\\hline
Indicadores ambientales y analítica & Media & Evidencia o beneficio incompletos; serie, asociación y corte verificable. & Conservar captura prioritaria; diferir cálculo/exportación y documentar atraso sin inventar resultado.\\\hline
Contenido institucional no transaccional & Baja & Información ordinaria diferida; no comprende portal de atención o avisos urgentes. & Mensaje de indisponibilidad y recuperación posterior; sin afectar canales críticos.\\\hline
\end{longtable}\endgroup
\FuenteTabla{Clasificación propuesta de Synaptix; Caso 07, caps.~3/15/19, y BA art.~78.}
Cada clase mantiene disponibilidad mensual mínima, respuesta y resolución: crítica 99,9\,\%, 15 min y 4 h; alta 99,5\,\%, 1 h y 8 h; media 99,0\,\%, 4 h y 24 h; baja 98,0\,\%, 8 h y 48 h. Crítica/alta disponen de atención 24\,$\times$\,7; media, horario hábil extendido; baja, horario hábil. Un incidente crítico conserva informe de causa raíz en cinco días hábiles y uno alto en diez. Disponibilidad se mide sobre la transacción de extremo a extremo; éxito de nube, réplica o pantalla no sustituye recepción y efecto de negocio. \parencite[art.~78]{sd8-ba}

\subsection{Desconexión, falla local y captura alternativa}\label{sec:sd8-bcp-local}
\textbf{BC-01, pérdida WAN (R-02/R-03/R-08).} SRE distingue enlace de energía y nodo. Se retiran los dos enlaces en el ensayo; la Célula Local mantiene recepción, asignación, salida/regreso, legajo, combustible y escuela durante 24 horas, con permisos locales autorizados y diario durable. Pantalán, varadero y escuela conservan 12 horas en terminales; Ilque, ocho horas sin cobertura. Se advierte antigüedad de copias y funciones externas pendientes. No se suponen 24 horas de batería de cada dispositivo ni combustible de generador por perder Internet. Energía y redundancia mantienen el diseño dimensionado de SD4, sin ampliar su autonomía por esta declaración.

En mostrador se identifica embarcación/persona, se verifica documentación local vigente, se prepara legajo en no más de 90 segundos y se registra el hecho con ID, secuencia, operador, instante y resultado. Si la autoridad requiere un canal indisponible, el legajo queda preparado y su remisión se realiza por el medio autorizado disponible; la plataforma no otorga zarpe ni interpreta silencio como autorización. No están disponibles confirmación de pagos/saldos externos, consulta actual a una autoridad caída, tratamiento remoto de IN3, federación remota ni analítica central fresca; se difiere el efecto y se muestra su estado, se usa captura/revisión manual de póliza cuando proceda y se conserva la tarea. El permiso local vigente no se sustituye por una contraseña compartida ni por una autenticación remota inexistente. \parencites[RT-03.10/13; RT-09.01; numeral 16.1, decisión 21]{sd8-c07}[RT-03.13]{sd8-bt}

Al reconectar, Sincronización envía automáticamente hechos desde su diario, verifica huella/secuencia e idempotencia, aplica autoridad por dominio y concilia sin pérdida. La prueba exige que después de 24 h de desconexión termine en 30 minutos como máximo, sin intervención manual en esa sincronización y sin perder salida, regreso, despacho o asistencia. Un error conserva el original y abre incidente; no se resuelve con última escritura que borre un hecho de seguridad, ni se comunica éxito porque la cola esté vacía. Un registro alternativo en papel es otro escenario y exige incorporación controlada; no cuenta como prueba de operación electrónica desconectada.

\textbf{BC-02, falla de nodo local (R-02/R-13).} SRE detecta mediante salud de servicio y transacción, aísla el nodo que falla y promueve únicamente la réplica local que conserve corte y exclusión del escritor anterior. Verifica energía, red, diario, permisos y captura completa antes de declarar servicio. Si falla también la alternativa o no puede excluir escritura doble, se bloquea la función afectada, se preserva almacenamiento y se activa respaldo manual seguro; no se confunde esa condición con simple pérdida WAN ni se exige administración técnica a la marina. Reparación y resiembra son tareas Synaptix; el titular de proceso valida negocio.

\textbf{BC-03, incertidumbre sobre alumnos, custodia o navegación (R-10/R-11).} El referente de escuela y su suplente contrastan lista local, salida real, acompañante, autorización y retirada. Ante discrepancia se mantiene vigilancia y se activa la respuesta humana del proceso; no se inventa presencia a partir de matrícula ni ausencia de mensaje. Operaciones contrasta salida/regreso y custodia por radio y contacto, sin concluir emergencia ni regreso automáticamente por telemetría tardía. Sin autorización probatoria no se entrega un menor o embarcación. La comunicación no exige tocar pantallas durante amarre o izaje; se realiza antes/después en condición segura.

\textbf{BC-04, aviso masivo fallido o marejada (R-01/R-13).} Notificaciones conserva destinatario, intento, acuse técnico y confirmación humana separados. Sin acuse, el operador abre llamada con titular/suplente, instrucción, hora, respuesta y siguiente escalamiento. Radio y teléfono disponible sirven de alternativa para expuestos, con evidencia de intento y contacto efectivo. Aviso oficial de marejada suspende toda actividad programada del proyecto; la respuesta a una emergencia real conserva su procedimiento. Se comunica al CLIENTE afectación, medidas y próxima actualización; no se utiliza continuidad para seguir una faena prohibida.

\textbf{BC-05, ausencia de conocedor del legado o cierre fallido (R-04/R-12).} Datos usa copia/inventario y procedimiento de cierre transferidos desde mes 6, con suplente técnico; Administración usa el suplente funcional ensayado. Se conserva último corte comprobado, se verifica recuento/saldo y se mantiene convivencia mientras no haya evidencia suficiente. No se calcula deuda por inferencia ni se reemplaza interpretación con un respaldo sin probar. La conciliación diaria y cierre mensual durante marcha blanca mantienen las puertas de SD5; un dato no explicado bloquea el corte.

Todo registro alternativo usa folio único, proceso, hecho, operador, hora, causa y relación con autorización; contiene sólo los datos imprescindibles. El titular funcional custodia físicamente la hoja bajo acceso restringido y SRE/Datos incorporan una imagen y el hecho trazable al ID original, con revisión independiente de recuento, efectos y duplicados. Se separa captura de validación: incorporar un papel no permite crear un segundo cargo, salida o despacho. El retorno requiere conciliar todos los folios, resolver pendientes y conservar evidencia por la retención del dominio de SD5; no se reinstala un cuaderno como sistema oficial permanente. La contingencia manual reduce daño pero no acredita la disponibilidad electrónica ni subsana una brecha del diseño.


\section{Plan de Acción a Riesgos}
Cada ficha T-16 selecciona prevención, mitigación y contingencia con titular, disparador, plazo, evidencia y paquete. Proyecto comprueba capacidad y cuenta de costo antes de autorizar. La mitigación reduce exposición; la contingencia se activa por el evento; un compromiso de control no acredita eficacia ejecutada.

\subsection{Selección de respuestas y costo-beneficio}
La comparación se expresa con esfuerzo, capacidad, efecto en plazo y resultados verificables; importes y valorización se conservan en la Oferta Económica. Se priorizan respuestas que cumplen cláusulas y reducen retrabajo o indisponibilidad sin crear una segunda imputación. \parencite[art.~50.2, p.~29]{sd8-ba}

El control temprano de variante y portabilidad se integra en selección e IaC; su beneficio es evitar una recepción incompatible y permitir traslado íntegro. Los dos ensayos y pruebas negativas consumen los paquetes de calidad/migración ya declarados, y su beneficio se comprueba en conciliación y retorno. La formación por perfil utiliza cohortes y acompañamiento de SD7 para impedir habilitación sin competencia. Las ventanas alternativas y el refuerzo propio sólo son una respuesta suficiente si su red, disponibilidad y costo están demostrados. La seguridad de personas y la protección de menores no se subordinan a una razón monetaria de conveniencia.

Las acciones iniciales de gobierno se imputan a 1.1.2.1, 32 HH de registro y respuestas; la revisión utiliza 1.1.4 y el servicio mensual correspondiente. Esas horas son trabajo comprometido, no una reserva libre. Cada titular informa actividad, horas efectivas y evidencia; Proyecto concilia el consumo con T-15 y la cuenta económica, sin publicar tarifas en este documento.

\subsection{Contingencia, gestión y reflejo en cronograma}
Se conserva la contingencia marina de diez días hábiles por campaña y las dos ventanas alternativas de SD7, consumidas una sola vez y antes del hito. No se suman automáticamente al escenario de seis días de repetición ni a los déficits CPM/PERT. La ejecución que sobrepase el escenario exige capacidad propia adicional comprobada, sin reducir calidad ni mover meses contractuales.

Se selecciona reserva de gestión temporal de cero días adicionales sobre los hitos obligatorios. Tampoco se presenta una bolsa de HH libres sin personal y agenda. Los riesgos no identificados se financian con gestión interna del precio fijo y refuerzo de Synaptix, con autorización del Ejecutivo y cálculo de esfuerzo/costo antes de ejecutar; su valorización y provisión permanecen en la Oferta Económica. Esto no permite cambiar objeto, garantías o calendario, ni cobrar un incumplimiento como ampliación de alcance. \parencite[art.~72, p.~38]{sd8-ba}

El registro de consumo conserva riesgo, decisión, reserva de origen, unidad, saldo inicial, cantidad usada, saldo final, actividad, ejecutor y cuenta. Proyecto comprueba $S_f=S_i-C$ y que un mismo consumo no aparece en dos riesgos. T-15 recibe cada acción con predecesora, ventana y responsable; si no cabe, se bloquea la selección y se dimensiona refuerzo. Holgura CPM, calendario protegido y tiempo de retorno no se tratan como la misma reserva.

\subsection{Acciones, ejecutores y carga de preparación}
La acción AC-ID mantiene el mismo riesgo ID de T-16 y hereda mitigación, contingencia, disparador, plazo, paquete y evidencia de su ficha; el registro \path{componentes/acciones_riesgos_control.csv} añade ejecutor y relevo por perfil. La Tabla~\ref{tab:sd8-ejecucion} concreta quién produce la respuesta y quién debe sustituirlo. Datos ejecuta corrección/carga de R-12 y Calidad verifica de forma independiente; propietario no significa ejecutor de todas las tareas. Una suplencia por perfil no acredita una persona disponible ni competencias certificadas.
\begingroup\setlength{\tabcolsep}{3pt}
\begin{longtable}{L{.13\textwidth}L{.31\textwidth}L{.34\textwidth}L{.13\textwidth}}
\caption{Asignación propuesta y carga inicial de control}\label{tab:sd8-ejecucion}\\\hline
\EncabezadoTabla{Riesgo} & \EncabezadoTabla{Ejecutor propuesto} & \EncabezadoTabla{Suplencia requerida} & \EncabezadoTabla{Carga inicial}\\\hline\endfirsthead
\hline\EncabezadoTabla{Riesgo} & \EncabezadoTabla{Ejecutor propuesto} & \EncabezadoTabla{Suplencia requerida} & \EncabezadoTabla{Carga inicial}\\\hline\endhead
R-01 & José Lara Arce & refuerzo del perfil O/L T-15 & 12 HH\\\hline
R-02 & José Lara Arce & refuerzo del perfil O/L T-15 & 12 HH\\\hline
R-03 & Fernando Guerrero & refuerzo de Arquitectura/Integración T-15 & 12 HH\\\hline
R-04 & Ignacio Silva & refuerzo del perfil J T-15 & 12 HH\\\hline
R-05 & Ignacio Silva & refuerzo del perfil J T-15 & 8 HH\\\hline
R-06 & Fernando Guerrero & refuerzo de Arquitectura/Integración T-15 & 8 HH\\\hline
R-07 & Fernando Guerrero & refuerzo de Arquitectura/Integración T-15 & 8 HH\\\hline
R-08 & Fernando Guerrero & refuerzo de Arquitectura/Integración T-15 & 8 HH\\\hline
R-09 & Bruno Toro & refuerzo de Seguridad T-15 & 12 HH\\\hline
R-10 & Bruno Toro & refuerzo de Seguridad T-15 & 12 HH\\\hline
R-11 & Fernando Guerrero & refuerzo de Arquitectura/Integración T-15 & 12 HH\\\hline
R-12 & Nicolás Torfan en corrección/carga; Davor Serey en verificación independiente & refuerzos D y Q de T-15, distintos para producir y verificar & 12 HH\\\hline
R-13 & José Lara Arce & refuerzo del perfil O/L T-15 & 8 HH\\\hline
R-14 & José Lara Arce & refuerzo del perfil O/L T-15 & 8 HH\\\hline
R-15 & Ignacio Silva & refuerzo del perfil J T-15 & 8 HH\\\hline
R-16 & Ignacio Silva & refuerzo del perfil J T-15 & 8 HH\\\hline
R-IN1-01 & José Lara Arce & refuerzo del perfil O/L T-15 & 8 HH\\\hline
R-IN1-02 & José Lara Arce & refuerzo del perfil O/L T-15 & 8 HH\\\hline
R-IN1-03 & Fernando Guerrero & refuerzo de Arquitectura/Integración T-15 & 8 HH\\\hline
R-IN1-04 & Bruno Toro & refuerzo de Seguridad T-15 & 12 HH\\\hline
R-IN2-01 & José Lara Arce & refuerzo del perfil O/L T-15 & 8 HH\\\hline
R-IN2-02 & José Lara Arce & refuerzo del perfil O/L T-15 & 8 HH\\\hline
R-IN2-03 & Ignacio Silva & refuerzo del perfil J T-15 & 12 HH\\\hline
R-IN2-04 & Fernando Guerrero & refuerzo de Arquitectura/Integración T-15 & 8 HH\\\hline
R-IN3-01 & Fernando Guerrero & refuerzo de Arquitectura/Integración T-15 & 8 HH\\\hline
R-IN3-02 & José Lara Arce & refuerzo del perfil O/L T-15 & 8 HH\\\hline
R-IN3-03 & Bruno Toro & refuerzo de Seguridad T-15 & 12 HH\\\hline
R-IN3-04 & José Lara Arce & refuerzo del perfil O/L T-15 & 8 HH\\\hline
R-IN3-05 & Ignacio Silva & refuerzo del perfil J T-15 & 8 HH\\\hline
R-IN4-01 & Ignacio Silva & refuerzo del perfil J T-15 & 12 HH\\\hline
R-IN4-02 & Ignacio Silva & refuerzo del perfil J T-15 & 12 HH\\\hline
R-IN5-01 & Ignacio Silva & refuerzo del perfil J T-15 & 8 HH\\\hline
R-IN5-02 & José Lara Arce & refuerzo del perfil O/L T-15 & 8 HH\\\hline
R-IN5-03 & Bruno Toro & refuerzo de Seguridad T-15 & 12 HH\\\hline
\end{longtable}\endgroup
\FuenteTabla{Titulares y perfiles de T-16/T-15; asignaciones y estimaciones propuestas por Synaptix.}
Cada preparación presupone cuatro HH para expediente y decisión, cuatro para revisión técnica y otras cuatro para guion de contingencia si $I=5$. El total es 328 HH de trabajo de control, a distribuir por perfil y antes del plazo de la ficha. No mide ensayos integrales, corrección de software, obras ni cobertura; se sustituye por la carga real de sus actividades antes de autorizar. Las 32 HH de registro de 1.1.2.1 se concilian con estos resultados: no se suman automáticamente si ya financian la misma tarea. El complemento identifica necesidad de recursos, no una segunda reserva de ejecución.

Proyecto selecciona una acción sólo con ID de persona o puesto, perfil, competencia, permisos, dedicación, jornada, ausencias, suplente y agenda. Para cada período $C_p=J_p-A_p-G_p$ y $S_p=C_p-\sum H_{a,p}$; gobierno, guardia realmente utilizada y otros paquetes forman parte de la carga, no del saldo libre. Las funciones SRE/Implantación de José usan una sola identidad y dos actividades coincidentes requieren otro ejecutor. El control del T-15 vigente declara HH y agendas globales incompletas: no se considera demostrada capacidad de las 34 acciones mientras falten esas entradas. Proyecto mantiene el bloqueo y dimensiona refuerzo del perfil ya previsto, sin inventar una nómina contratada.

\subsection{Comparación costo-beneficio y selección}
El costo técnico comprende HH de preparación más trabajo EDT, capacidad mantenida, ventanas y complejidad de recuperación. El beneficio comprende obligaciones preservadas y carga o demora evitada. Se compara la alternativa real y no se monetiza $P I$; OE contiene precios, tarifas y provisiones. La Tabla~\ref{tab:sd8-beneficio} selecciona respuesta por mecanismo.
\begingroup\setlength{\tabcolsep}{3pt}
\begin{longtable}{L{.13\textwidth}L{.29\textwidth}L{.30\textwidth}L{.19\textwidth}}
\caption{Costo técnico, alternativa y beneficio esperado}\label{tab:sd8-beneficio}\\\hline
\EncabezadoTabla{Riesgos} & \EncabezadoTabla{Carga de respuesta} & \EncabezadoTabla{Alternativa y efecto} & \EncabezadoTabla{Selección}\\\hline\endfirsthead
\hline\EncabezadoTabla{Riesgos} & \EncabezadoTabla{Carga de respuesta} & \EncabezadoTabla{Alternativa y efecto} & \EncabezadoTabla{Selección}\\\hline\endhead
R-02/R-08 & Pasivo preparado: 18 vCPU/72 GiB; ampliar al peak y probar antes de abrir. & Réplica permite conmutación; respaldo solamente requiere 233,33 min para 60 GB bajo hipótesis. & Elegir pasivo y restauración independiente; no sustituir RTO por disponibilidad de un archivo.\\\hline
R-03/R-12 & Dos ensayos, conciliación y retorno; Datos produce y Calidad revisa. & Un ensayo único consume menos preparación pero no satisface la obligación; 6 días adicionales sólo modelan repetición parcial. & Elegir ambos ensayos y reservar corrección del defecto; no recortar evidencia para ahorrar.\\\hline
R-04/R-14 y adopción & Preparación de referentes/suplentes y evaluación de tarea, más cohortes por perfil. & Formación genérica reduce trabajo aparente pero no acredita cierre, permisos ni uso correcto. & Elegir práctica por proceso, recuperación y habilitación individual; bloquear acceso no preparado.\\\hline
R-06/R-11 & Recepción de variante, montaje y compatibilidad antes de habilitar. & Compra sin revisión puede ahorrar control inicial pero obliga a sustitución y pierde la ventana admisible. & Evitar intervención insegura; sustituto verificado y nueva recepción.\\\hline
R-09/R-10 e IN3/IN5 sensibles & Revisión negativa, procedencia, permiso y minimización; intrusión independiente cuando corresponde. & Abrir permiso o usar datos reales para simplificar invalida protección y evidencia. & Reducir exposición con controles; no aceptar incumplimiento de privacidad.\\\hline
R-01/R-15 y ventanas IN2 & Diez días de escenario por campaña y dos alternativas, recursos por ventana. & Mover hito o continuar trabajo durante aviso infringe restricciones; HH sin ventana no recuperan calendario. & Suspender, recalcular red y dimensionar refuerzo propio después del cese.\\\hline
R-07/R-16 & Exportación, custodia y paquete reproducible en cada liberación. & No depositar reduce trabajo pero aumenta dependencia y hace imposible activar la salida. & Mantener titularidad y transferencia; soporte del proveedor no elimina responsabilidad Synaptix.\\\hline
IN1--IN5, eficacia de beneficio & Revisión de grupo, serie y universo de casos; registro de beneficio adverso. & Ocultar casos difíciles o proyectar ahorro sin población comparable aparenta éxito sin evidencia. & Verificar universo, conservar servicio y reevaluar o suspender incentivo no sustentado.\\\hline
\end{longtable}\endgroup
\FuenteTabla{Decisiones de Synaptix con SD4/SD5/T-15, T-16 y restricciones del Caso 07. Beneficios previstos, no observados.}
La estrategia es reducción para fallas de datos, seguridad, adopción y capacidad; evitar para maniobra insegura o promoción sin prueba; compartir provisión/custodia con el tercero competente sin transferir obligación del contrato; y aceptación sólo del residual evidenciado dentro de competencia. Ninguna respuesta permite cámaras sobre menores, interacción durante izaje o amarre, dispositivo prohibido o retención por mora sin decisión formal.

\subsection{Reservas: dimensionamiento, activación y saldo}
La reserva marina es un escenario de diez días hábiles por campaña, bajo una sola cuenta R-01/PL-R01. La contingencia PL-R03 es hasta seis días adicionales para repetición parcial; no equivale a las dos ejecuciones obligatorias. Se selecciona reserva de gestión de cero días adicionales sobre meses contractuales; su provisión económica y refuerzo propio requieren decisión del Ejecutivo y actividad/recursos definidos. No se adopta un porcentaje universal ni una bolsa de HH que no tenga agenda.

Para una corrección con $n$ personas y $h$ HH efectivas por día, $d$ días de respuesta requieren $n h d$ HH antes de añadir traslados o revisión independiente. Como hipótesis localizada, dos ejecutores a cinco HH/día durante seis días requieren 60 HH; durante diez días, 100 HH. No demuestran que esas personas estén disponibles ni que un ensayo completo quepa en seis días. Si las dos respuestas comparten ejecutores o son consecutivas, sus demandas suman 160 HH y 16 jornadas de equipo, no diez; si coinciden en tiempo, se calcula carga diaria y se aporta otro equipo, sin reutilizar la misma agenda.

El acta de activación contiene riesgo, causa, respuesta elegida, predecesoras, ventana válida, carga por ejecutor, recurso/suplente, reserva de origen y saldo. Proyecto propaga cada respuesta al cronograma antes del hito y verifica margen después de congelamientos y revisión del CLIENTE. La reserva sólo se consume tras aprobación competente; $S_f=S_i-C$ conserva unidad y no mezcla HH, días, capacidad de puesto o dinero. La falta de cabida se registra como brecha, no se corrige moviendo producción o reduciendo marcha blanca. La reserva de gestión no resta derechos de recepción. \parencites[arts.~17, 50.2 y 72]{sd8-ba}[RT-19.03/04]{sd8-bt}

\subsection{Respuesta técnica y plazos de seguridad}\label{sec:sd8-respuesta-seguridad}
Estas reglas complementan AC-R-09/10/13 y las acciones de IN3, conservando ejecutor, propietario y paquete de cada ficha. Seguridad dirige evaluación/contención; SRE ejecuta aislamiento y restitución autorizados; Datos preserva evidencia y concilia; Calidad verifica corrección y negocio. Proyecto coordina comunicación y capacidad. El catálogo de escenarios de 8.2 remite a SEC-1.0/IR-1.0, SEC-ROT-1.0/GLS-1.0, DEP-6, AUD-5 e IN3-C01--C05; no se declara eficacia por escribir estas instrucciones.

El programa de vulnerabilidades mantiene activo, hallazgo, severidad, fecha de publicación/detección, responsable, vencimiento, corrección y prueba. El máximo de remediación es siete días corridos para críticas, quince para altas y treinta para medias, desde publicación o detección, sin reiniciar reloj al asignar técnico ni confundirlo con SLA de incidente. Si no existe parche, se retira, sustituye o deshabilita la función afectada y se comprueba remediación; control compensatorio sin evidencia o congelamiento no extienden el plazo. Se aplica regresión y retorno antes de reabrir, conservando hechos y evidencia. \parencites[art.~21.1]{sd8-ba}[RT-11.04]{sd8-bt}

IR-1.0 abre cronología desde detección y escala inmediatamente el crítico. El CLIENTE recibe comunicación de incidente crítico dentro de dos horas desde detección; toda brecha de seguridad o datos personales se notifica dentro de veinticuatro horas con informe preliminar y análisis de causa raíz dentro de los cinco días hábiles siguientes. Un incidente crítico que además sea brecha cumple ambos avisos: no espera al de veinticuatro horas para el primero. Seguridad/SRE preservan original, alcance conocido, incertidumbre, medidas, responsables y próxima actualización; Proyecto coordina los canales funcionales con la Contraparte, sin esperar comité ni cierre de investigación. \parencites[art.~21.3]{sd8-ba}[RT-11.18/19]{sd8-bt}

Ransomware o fuga requieren aislar el activo/frontera afectado, revocar credencial, custodiar evidencia y comprobar copias/claves antes de restaurar; la captura local segura se preserva sin aislar simultáneamente el par por defecto. DDoS o saturación limitan tráfico diferible y aplican capacidad/prioridad de SD4 sin descartar registros críticos. Dispositivo comprometido pierde facultades y sincroniza hechos sólo tras comprobación; una orden remota durante WAN caída no se presenta como ejecutada. Pérdida de observabilidad mantiene incidente abierto y escalamiento alternativo hasta comprobar detección y respuesta, sin declarar cobertura por contar un puesto SOC nominal.

La habilitación o rehabilitación de IN3 exige comprobar IN3-C01--C05 y el modelado propio: no escritura sin revisión, fragmento/versión trazables, permiso y tratamiento autorizado, resultados por campo/formato, desactivación e ingreso manual con conciliación de respuestas tardías. Región, versión/configuración y garantía de no entrenamiento permanecen puerta previa a tratar documentos del CLIENTE. Proyecto no consume una reserva distinta por cada escenario del mismo incidente; recursos, pruebas y suplencias siguen sujetos a la conciliación de T-15.


\subsection{Gobierno de activación y recuperación}\label{sec:sd8-gobierno-recuperacion}
Ante R-02/R-03/R-12/R-13, SRE y Seguridad evalúan el incidente de su ámbito y activan el procedimiento autorizado; Proyecto comunica y coordina, y el referente funcional decide la suspensión segura de la actividad afectada. La pérdida de visibilidad de un menor o de integridad de custodia exige respuesta inmediata y no espera una puntuación. Datos preserva los originales y coordina conciliación; Calidad comprueba evidencia y el CLIENTE conserva aceptación funcional. El retorno a operación normal requiere comprobar servicio, integridad y registros pendientes antes de levantar el bloqueo. Los procedimientos manuales, secuencias, tiempos y ensayos se desarrollan en las secciones~\ref{sec:sd8-bcp-bia} a~\ref{sec:sd8-pruebas-continuidad}; ninguna función aquí asignada acredita suplencias ni ensayos realizados.

\subsection{DRP: conmutación regional, retorno y custodia}\label{sec:sd8-drp}
El par seleccionado es AWS São Paulo, \texttt{sa-east-1}, y AWS México Central, \texttt{mx-central-1}, activo-pasivo preparado según DR-4 de SD4. La separación orientativa es 7.613,84 km entre referencias urbanas, no centros de datos ni latencia. Reduce amenazas físicas locales; no elimina dependencia común de AWS, organización/IAM, DNS, federación, claves, despliegue, corrupción o rutas de enlace. SRE comprueba esas dependencias y Seguridad autorizaciones de tratamiento Brasil/México y custodia; un traslado regional no acredita permiso legal por sí mismo. No se mantiene Google Cloud Santiago como un segundo destino simultáneo ni México como región temporal.

\textbf{DR-01, detección y decisión.} La prueba mide RTO desde la interrupción y RPO desde el último corte común recuperable; SRE abre incidente y comunica al CLIENTE, Seguridad analiza integridad y Proyecto coordina consecuencias. Las alarmas de SD4 usan latido SQL y objeto, manifiestos y observación cada minuto: más de 120 s en dos muestras avisa; más de 300 s en dos genera incidencia; más de 600 s en una escala y suspende liberaciones/cargas no críticas; 900 s, fallo o ausencia de dos muestras válidas bloquean declaración conforme de RPO. Se contempla reloj no fiable y objetos pendientes; \texttt{ReplicaLag} no demuestra por sí solo copia de toda la evidencia. El incumplimiento del RPO se registra sin alterar su mínimo.

\textbf{DR-02, exclusión y corte.} Arquitectura/SRE ejecutan el procedimiento de aislamiento, detienen emisores y escrituras del origen y conservan ID/época de autoridad. Antes de promover, Datos verifica corte común de los dos almacenes, objetos, permisos, revocaciones, eliminaciones y bandejas de hechos. Una partición que impida comprobar exclusión del escritor anterior no autoriza promoción; se mantiene operación local segura y se escala. DNS por sí solo no impide a un escritor antiguo producir efectos. Si la evidencia del diseño no resuelve exclusión, identidad de emergencia o clave, queda una brecha de arquitectura y no se declara alcanzado el RTO por saltar el control.

\textbf{DR-03, preparación y restitución.} SRE promueve réplicas según SD4, ajusta capacidad al perfil productivo vigente, activa infraestructura versionada y comprueba cuotas, certificados, claves, identidad regional y ruta de acceso desde cliente y recinto. No se replica una contraseña ni un secreto externo por suponer que el proveedor lo haga. El pool de contingencia conserva identidad estable y mecanismos de recuperación preparados de SD4; federación caída no se usa como único acceso para respuesta. Seguridad comprueba permisos mínimos y autenticación de responsables de emergencia. Colas se reconstruyen desde inbox/efecto/outbox durables: los servicios de transporte no son la única copia. Se consultan respuestas externas inciertas por ID antes de repetir, y se bloquea un efecto sin idempotencia/consulta hasta conciliación competente.

Calidad y referente funcional prueban recepción/amarra, salida/regreso, presencia/retirada escolar, combustible, avisos, integridad y consulta/efecto administrativo crítico. Se reanudan consumidores y se abre tráfico sólo tras esas comprobaciones. SRE comunica restablecimiento y restricciones restantes; la duración del incidente continúa hasta recuperación efectiva del negocio, no hasta que un proceso arranque. La Contraparte recibe procedimiento y entrenamiento para solicitar, autorizar o coordinar la conmutación desde la consola gestionada; la ejecución técnica pertenece al servicio Synaptix, sin exigir un administrador AWS local inexistente. \parencites[RT-07.01--07]{sd8-bt}[numerales 2.4 y 17.6]{sd8-c07}

\textbf{Presupuesto técnico de RTO, hipótesis por validar.} Se asignan 15 min a detección/decisión, 30 a exclusión y corte común, 30 a promoción, 45 a capacidad/rutas/identidad y 90 a prueba integrada y reconciliación: 210 min; los 30 restantes son margen de ese escenario, no una nueva reserva de calendario contractual. Las fases solapadas no se descuentan sin prueba. Fallo de clave, identidad, datos o exclusión puede exceder esa secuencia: el plan conserva la brecha y corrige antes de producción. No se presenta el presupuesto como medición realizada ni garantía de capacidad faltante.

\textbf{DR-04, retorno.} Se estabiliza el primario sin servir escritura, se realiza copia/replicación desde la autoridad actualmente activa en México, se reconcilian hechos, objetos, permisos y efectos externos, y se ensaya el retorno de los pendientes de borde. Datos verifica corte y Calidad negocio; SRE ejecuta un cambio de autoridad autorizado con aislamiento del escritor anterior y observación reforzada. Se mantiene origen de contingencia protegido hasta confirmar eficacia y aceptación; no se reutiliza una réplica divergente ni se eliminan evidencias para aparentar normalidad. Proyecto registra residual observado, duración, último dato recuperable, condiciones de retorno y acciones correctivas.

\subsection{Respaldo, restauración y conservación}\label{sec:sd8-backup}
Se aplica 3-2-1-1-0: tres copias recuperables del conjunto, dos medios distintos, una fuera de sitio, una inmutable o fuera de línea y cero errores en la verificación. Datos/SRE mantienen producción/réplica, respaldo versionado y custodia externa independiente prevista en SD4, incluyendo original, documentos y diario local pendiente. Una réplica que propaga corrupción no sustituye la copia de recuperación; tres referencias al mismo objeto no son tres copias. La política requiere inventario de medio, ubicación, propietario, corte, huella, cifrado, clave, retención, expiración y resultado de prueba. La variante de custodia y sus recursos se verifican antes de confiarle recuperación, sin afirmar contrato ya firmado.

Claves de respaldo se custodian de forma independiente del conjunto respaldado, con acceso nominal, doble control para recuperación de emergencia y prueba de descifrado desde el entorno de restauración. Copias inmutables tienen regla de retención que impide borrado/modificación aun con credenciales administrativas comprometidas; Seguridad ensaya el intento denegado. Una pérdida de clave, copia incompleta o restauración con datos vencidos impide abrir tráfico. Se restaura aislado y se reaplican revocaciones/eliminaciones autorizadas y bloqueos probatorios antes de servir. \parencites[art.~20]{sd8-ba}[RT-07.09--13]{sd8-bt}

Frecuencia propuesta para los doce dominios de SD5: copia diaria, recuperación a punto en el tiempo de 35 días en nube y doce copias mensuales; subconjunto local diario de treinta días. La tabla de tiempos de SD5 reúne nueve grupos de recuperación: Clientes/contratos, Operación marítima, Escuela, Varadero/contratistas, Combustible, Ambiente/consumos, Finanzas/cobranza, Infraestructura y Auditoría/seguridad. Sus tiempos de restauración individual estimados por SD5 son, en ese orden, 74, 54, 71, 65, 47, 54, 71, 51 y 47 minutos para su volumen histórico provisional. Son escenarios por grupo y no se suman para restauración conjunta, estimada en 233,33 minutos con los supuestos explicados en 8.2. Antes de aceptar RT-07.13, Datos desglosa esa estimación por NAV, AMA, ESC, CON, BOY, VAR, CTR, AMB, FAC, ADM, AUT y ALE, incluyendo avisos y sus acuses. Los nueve grupos no acreditan todavía tiempo completo por cada dominio; se conservan como cálculo parcial y se agregan subconjuntos y procesos nuevos al inventario antes de respaldarlos; los documentos de evidencia mantienen retención propia, que puede superar treinta días, 35 días o doce meses. No se confunde rotación de respaldo con retención legal de fuente.

SRE verifica copia diaria, corte y huella; mensualmente restaura muestra representativa con claves, permisos, documentos y hechos pendientes, mide el tiempo efectivo y Calidad valida integridad. Un error invalida la prueba y abre remediación con responsable/fecha y repetición; cero errores es criterio de aceptación, no un resultado previamente obtenido. Los datos locales no confirmados no se borran por antigüedad y se incluyen en respaldo/replicación local antes de purgar; al reconectar se contrastan manifiestos de destino. Las muestras de menores se minimizan y protegen sin convertir datos de prueba en una divulgación.

\subsection{Pruebas, ventanas y criterio de recuperación}\label{sec:sd8-pruebas-continuidad}
Antes de cada producción E1/E2, Calidad coordina con SRE/Datos/Seguridad pruebas de recuperación regional, dos enlaces WAN retirados durante 24 h, terminales 12 h e Ilque 8 h, sincronización menor o igual a treinta minutos, nodo local primario y alternativa, mensajes perdidos/duplicados, caída de identidad, clave no disponible, copia corrupta, objeto sin réplica y eliminación/restauración. Se mide inicio/fin, RTO/RPO, corte de datos, recuentos y huellas, cola y latencia de negocio. Cada prueba conserva versión, carga representativa, ejecutor, verificador independiente cuando corresponde, permiso, incidencia y decisión. No se presenta su planificación como resultado ejecutado.

En operación se realizan al menos dos conmutaciones regionales reales al año, con informe RTO/RPO y corrección de brechas; resiliencia por instancia, zona, dependencia, latencia y disco al menos semestral; restauración representativa mensual; seguridad ofensiva por tercero independiente anual y antes de cada producción. SRE incorpora las pruebas en 1.12.4 y Calidad en paquetes de aceptación de SD7/SD9, evitando duplicar una ejecución. Se programan fuera de congelamiento y ventanas críticas con aviso de al menos diez días hábiles; la fecha se anticipa para conservar intervalo máximo, no se posterga por conveniencia. Una emergencia no requiere esperar ese aviso de mantenimiento. \parencites[arts.~20 y 21.3]{sd8-ba}[RT-07.07/12; RT-10.05--07; RT-11.20; cap.~20]{sd8-bt}

El CLIENTE valida el modo de operación y comunicación; el servicio gestionado ejecuta técnica. El retorno requiere recuperar flujos y accesos, conciliar hechos alternativos y pendientes, comprobar retenciones y restricciones, retirar doble autoridad y emitir comunicación de restablecimiento. Una condición impeditiva mantiene bloqueo incluso con puntuación residual baja. Los registros de prueba y eficacia alimentan T-16 y los comités; sólo entonces se propone residual observado o reapertura. Permanecen visibles las brechas de capacidad, calendario y diseño señaladas por SD4/SD5/T-15, sin trasladarlas a una aceptación tácita.


\subsection{Eficacia, residual y cierre}
Los objetivos residuales de T-16 dependen de ensayos, habilitaciones y participación todavía por comprobar. Hasta disponer de evidencia se mantiene la evaluación inicial y el estado de tratamiento comprometido. Calidad confirma integridad y cobertura; Seguridad, privacidad/procedencia; SRE, continuidad; la Contraparte conserva recepción.

Proyecto registra residual observado y aceptación competente, o nueva acción cuando falla la eficacia. Un riesgo alto/crítico abierto se escala con causa, exposición y recuperación; una condición impeditiva bloquea producción aunque el residual objetivo sea bajo. Al cierre de etapa y en la salida se entregan riesgos vigentes, responsables, reservas consumidas y obligaciones pendientes; no se elimina un antecedente para aparentar cumplimiento.
